\section{Cosmologia}\index{Cosmologia}\hypertarget{cosmologia}{}\label{cosmologia}

\begin{enfasi}{
E' più facile dominare su chi non crede in niente (La Storia Infinita, Kmorf)

\medskip

Tu credi che c'è un Dio solo? Fai bene; anche i demoni lo credono e tremano! (Giacomo Il Giusto 2, 19)}\end{enfasi}

\begin{narratore}[I Patroni]

In OBSS le divinità sono diverse dai tradizionali dei nei giochi di ruolo.

Le divinità, i Patroni, amano sporcarsi le mani, partecipare nelle faccende delle creature che le adorano, per loro è una sfida continua ad avere più credenti, adepti e persone più simili, per Tratti, a loro.

I Patroni sono stati creati come \href{https://www.treccani.it/vocabolario/parossismo/}{parossismo} dell'animo umano, dove tutto è un eccesso. Come spiriti liberati dal vaso di Pandora hanno il solo scopo di portare i loro Tratti al dominio rendendoli i più comuni e presenti tra le creature, specialmente tra le più potenti.
\end{narratore}

\medskip

\begin{multicols}{2}

In principio era il nulla che in sé conteneva il tutto.

L'Energia dei Patroni della Genesi, derivante dalle più primordiali pulsioni esplodeva in tutta la sua potenza senza alcun controllo

Manifestate come due lingue di una fiamma unica Ljust e Calicante sono la fonte infinita di infinita energia.

Ljust l'energia positiva, il calore, la luce, la vita e sintropia; Calicante, l'energia negativa, gelido odio, distruzione, morte ed entropia.

\textbf{Ljust} \index{Ljust}è la rappresentazione di ciò che luce ed vita portano sempre con se. Rappresenta la purezza del sentimento d'amore, la protezione della vita, il rispetto per l'altro, la curiosità per il nuovo, la voglia di migliorarsi sempre, la forza di combattere con coraggio e valore per il bene comune. E' la spinta vitale al cambiamento, il caos che evolve ma non distrugge.

\textbf{Calicante}\index{Calicante} è la rappresentazione del buio, dell'odio, della rabbia e violenza. Calicante è vendetta e fredda distruzione, non c'è interesse in alcuna forma di vita piuttosto le usa, le sfrutta e solo in certi casi ne subisce la presenza. Ama sadicamente la sofferenza. E' l'entropia che annienta e annichilisce e trova piacere nel farlo.

\textbf{Atmos} \index{Atmos}è il testimone, colui che segna il passaggio del tempo e trascrive ogni accadimento sulla Terra e trai Patroni della Genesi. Una entità nata dalla creazione perché prevenisse la distruzione assoluta. Vigila e trascrive quanto fanno i Patroni della Genesi, le divinità che hanno generato il creato.

Un quarto essere fu presente, il cui nome non ci è mai giunto, ma che conosciamo come \emph{Ominiessenza}, colui che era il primo nato dai Patroni della Genesi.

Il Primogenito dei Patroni della Genesi curiosando nell'universo tutto giunse sulla Terra e qui, non si sa come, venne catturato dalla Freten.
Catturato, fatto a pezzi, collegato a macchine per poter attingere alla sua eterna energia. Non si sa come è stato possibile e tuttora è uno dei più grandi misteri irrisolti.

Quando il primo reattore della Frenten venne collegato si aprì una \emph{frattura} e fu così che i Patroni della Genesi scoprirono cosa era accaduto e l'ira di Calicante fu immensa e totale.

La sua volontà e furia fu talmente soverchiante che Ljust non poté intervenire immediatamente, non poté impedire la morte indiscriminata di innocenti e la quasi distruzione di un interno mondo.

Cosa successe in quell'anno lo possiamo ancora vedere nelle macerie e distruzione rimaste sulla Terra.

Ljust riuscì infine a placare, con sommo sacrificio, Calicante e questo richiamò le sue manifestazioni.

La morte del primo Patrono aveva creato numerosi squarci, fratture, nella realtà, Portali che si aprivano in altri mondi, altri universi. Portali che continuavano ad aprirsi e chiudersi cambiando posizione di continuo. In ognuno di questi Portali c'era un pò del loro figlio.

Ljust e Calicante decisero di fare di questo nostro mondo, la Terra, il loro spazio di gioco e qui decisero di fare divertire la loro futura progenie, i Patroni.

Decisero di comune accordo di generare un Patrono che sovrintendesse a questi Portali. Crearono un Patrono che fosse capace di percepire, aprire e bloccarli Così venne creato \textbf{Lynx}, il Guardiano dei Portali.

Lynx \index{Lynx}sovrintende al vuoto cosmico, all'accesso ai Piani, ai portali che con l'avvicendarsi di caos e ordine, di bene e male, di luce e tenebra stanno sempre più creando fratture al confine esistente la la Terra e l'Oltre.

Lynx li percepisce, li sente, sa dove si stanno generando o spegnendo, con il passare del tempo infatti alcuni di questi Portali sono divenuti stabili e definitivi, altri invece continuano a generarsi casualmente e sempre in modo totalmente ignoto rimangono attivi o si esauriscono. Viaggiando di continuo nel non luogo Lynx chiude i portali più grandi ma per uno che ne chiude un altro si apre. Lynx ha privato le Liste di Magia da molti degli incantesimi che agiscono sui piani, per proteggere la Terra ed i futuri Patroni dalle minacce esterne.

Lynx è il custode e carceriere di tutta la Terra.

Atmos, preoccupato per l'equilibrio sulla Terra, incanalò le energie primordiali e divine dei Patroni della Genesi andando a creare un Patrono che potesse rivaleggiare e tenere a bada gli altri futuri Patroni.

Il primo creato da Atmos, con l'aiuto di Ljust, e l'intervento di Calicante, fu \textbf{Gradh}\index{Gradh}, Patrono dell'Umanità (e di tutte le razze senzienti), colui che avrebbe difeso la Terra dalle creature esterne e dagli altri Patroni. Gradh racchiude in sé il dualismo dei due Patroni della Genesi, l'istinto innato alla protezione, alla difesa ed alla cura propri di Ljust e l'istinto di vendetta, violenza e furia di Calicante.

Si getta con coraggio nelle battaglie, attacca il nemico senza paura, protegge il più debole, difende la vita ma non teme di percorrere la strada della vendetta più distruttiva verso chi sfrutta e distrugge vite senza motivo. Gradh ama calarsi fra la gente e vivere con loro, come loro. Non si sente totalmente a suo agio nel pantheon con gli altri Patroni ne fra la gente comune, lui è Umano tra i Patroni e Patrono tra gli Umani. Passionale e gentile è il Patrono che maggiormente ha a cuore le sorti del nostro pianeta e delle sue creature. Gradh è il Patrono che ci protegge dai Patroni, forse l'unico vero figlio della Terra e dei Patroni della Genesi.

Le lingue di energie divine erano troppo intense, caotiche e pure perché Atmos potesse governarle per plasmare da solo i Patroni. Usando il grezzo potere dei Patroni della Genesi creò altri Patroni, la seconda venuta, ognuno influenzato in maniera diversa da Calicante o Ljust. Questi Patroni risultarono meno perfetti e divini rispetto alle sue intenzioni, più imperfetti ed \emph{umani} in quanto originati dalle emozioni, incontrollabili e pure dei Patroni della Genesi. Questi nuovi Patroni plasmano le volontà, fondano regni, comandano nell'ombra le pedine che osano chiedere i loro favori.

In questa apparente calma i Patroni perpetuano i loro interessi, diventare i più forti, i più importanti, coloro che hanno maggiori seguaci. Lo scopo è uno solo: avere più persone possibili che seguono i loro Tratti.

Se un Patrono agisce in prima persona o in modo indiscriminato sa che scatenerà la reazione di Gradh o l'intervento di Atmos che gli impediranno un uso incontrollato e massivo dei suoi poteri direttamente sulla Terra, come fecero i primissimi Patroni. Questo raramente li ferma e la stessa natura e creature tutte vengono spesso influenzate dal volere dei Patroni.

Ed è così che nascono sempre più spesso aberrazioni, malattie sempre nuove, terre maledette dove non può crescere nulla, per non parlare di pazzie che spesso coinvolgono chi invece dovrebbe proteggere i cittadini.
E' una dura vita quella dell'uomo comune che continuamente deve affrontare siccità o alluvioni, morie di animali ed un meteo irregolare se non assurdo, orde di creature venute dal nulla che vogliono solo sterminare tutti,

Ad ogni passo deve guardarsi intorno perché non può mai sapere chi ha venduto l'anima ad un Patrono per vivere un giorno in più.

Ovunque i nemici più forti sono i Draghi, generati e convocati nella prima venuta da Tàhil, il primo servitore di Calicante, unico, forse, dei primi Patroni ad essere rimasto.
Fanno incursioni al solo scopo di portare distruzione e morte, seminare paura ed orrore.

Molte regioni stanno diventando incubatrici di razze oscure e malvagie, legioni di non morti guidate da potenti negromanti si ammassano sui confini, i Draghi addestrano i loro corrotti adepti e scure spire nel cielo promettono tempesta.

Come la Freten riuscì a catturare il Primo figlio ed a soggiogarlo è tuttora un mistero.

\subsection{Patroni}\index{Patroni}\hypertarget{patroni}{}\label{patroni}

\begin{enfasi}{
Conan: A quali dei preghi?

Subotai: Io prego ai quattro venti e tu?

Conan: Io prego Crom, ma solo raramente... lui non ascolta. (Conan il Barbaro, film 1982)

\medskip

Infatti, come il corpo senza lo spirito è morto, così anche la fede senza le opere è morta. (Giacomo Il Giusto 2, 26. NdA Riferendosi ai punteggi dei Tratti collegati al Patrono...)
}
\end{enfasi}

Le creature tutte, anche chi non usa la magia possono sentire l'influenza di questi Poteri, di questi Patroni.

Ogni personaggio per il suo modo di essere (giocare) e comportarsi ha almeno un Tratto in comune con un Patrono e nel corso delle avventure e sua evoluzione matura e potenzia queste convinzioni potrà sentire maggiormente l'influsso ed effetti di un Patrono.

Non è necessario che abbia giurato fedeltà ad un Patrono o che ne sia Seguace o Devoto, sentirà comunque l'influenza del Patrono e riceverà dei doni da esso.

Un Patrono è ben contento se le creature seguono i suoi dettami, Tratti, e dona a coloro che lo fanno dei piccoli poteri come riconoscimento per la fedeltà a lui riservata, volutamente o meno. I poteri indicati sotto \emph{Tratti in Comune} sono cumulativi. Se non indicato diversamente i poteri sono usabili 1 volta al giorno e costano 2 Azioni.
Quando è indicato un incantesimo questo viene manifestato senza Prove di Magia o penalità dovute ad armatura.

Ogni \textbf{Patrono predilige uno o più forme energetiche}, se sei Seguace puoi usare quell'energia nelle tue magie, se sei un Devoto invece i tuoi incantesimi useranno una della forme energetiche indicate. Sono indicate le \hyperlink{listeprivilegiate}{\emph{Liste Privilegiate}} (pag. \pageref{listeprivilegiate})\index{Liste Magia privilegiate}, ovvero liste nelle quale il Devoto ha dei vantaggi d'uso.

Le forme di Energia vengono distinte tra fonti positive, neutrali e negative, servono per inquadrare meglio il vostro padrone, pardon, il Patrono che servite.

Fate la somma delle energie, se positiva il Patrono si può considerare buono, se a valore zero il Patrono è neutrale, se a valore negativo il Patrono è malvagio, per comodità nella lista delle forme di energia è indicato se il Patrono è \textbf{B}uono, \textbf{N}eutrale o \textbf{M}alvagio.

Nella descrizione del Patrono troverete anche la sua \textbf{manifestazione}, ovvero cosa accade quando un personaggio agisce in maniera particolarmente e significativamente consona ai Tratti seguiti dal Patrono. L'effetto è puramente scenico e di circostanza ma lascia sempre colpito chiunque lo possa osservare, solitamente garantisce un punto avanzamento in un Tratto collegato al Patrono. Non è necessario essere Devoto o Seguace, basta aver seguito in maniera particolarmente \emph{epica} quel Tratto.

E' presente anche l'indicazione dell'\textbf{arma preferita} dal Patrono. Non ci sono vantaggi nell'usarla, a meno di avere la specifica \hyperlink{Il Patrono è la mia Arma}{Abilità} (pag. \pageref{Il Patrono è la mia Arma}), la scelta è puramente personale e lasciata alla devozione del personaggio.

Sotto l'indicazione dell'arma preferita c'è l'indicazione della Regola \index{Regola} ovvero il comportamento che il Devoto deve cercare di rispettare.

Un incantatore che si affida ad un Patrono, almeno \textbf{2 Tratti} in comune, diventa un \textbf{Devoto}. Se ha almeno \textbf{1 Tratto} in comune e si affida ad un Patrono allora si dice che è un \textbf{Seguace}. Il \textbf{Vantaggio} indicato è solo per il Devoto.

\begin{giocatore}[Devoti e Seguaci]
Essere Devoti o Seguaci è una scelta vostra, nessuno ve la impone. Dovete sentirla come una occasione di gioco di ruolo, come un arricchimento del personaggio e non una costrizione. Essere Devoti o Seguaci non significa essere proni al volere del Patrono, anzi, significa essere ancora più convinto di propri Tratti, della propria personalità. \textbf{Un Patrono non chiede preghiere, ma chiede di essere se stessi}.
\end{giocatore}

Il personaggio potrebbe anche non seguire alcun Patrono pur avendo più Tratti in comune oppure potrebbe essere un Devoto o Seguace non del Patrono con cui hai più Tratti in comune od i Tratti a punteggio più alto. La scelta è sempre del personaggio e della sua sensibilità.

Le capacità acquisite legate ai Tratti in comune sono indipendenti dall'essere un Devoto, Seguace o semplicemente ateo, rappresentano i doni del Patrono a chi segue i suoi Tratti.

Nulla vieta che un personaggio riceva più poteri da Patroni diversi! Ad alti livelli quando il personaggio ha un alto punteggio nei vari Tratti posseduti questo capiterà frequentemente. \index{Vantaggi}

\begin{narratore}[Adattarsi]
Il Narratore può comunque concedere l'essere Seguace o Devoto pur se i Tratti non collimano perfettamente. Su richiesta del giocatore ed a sua discrezione può valutare la somiglianza di alcuni Tratti del personaggio a quelli del Patrono e valutarli idonei per esserne un Seguace o Devoto. In queste situazioni è necessario comprendere come il giocatore inquadra il personaggio e capire non solo se i Tratti ma anche il sentimento del personaggio è affine al Patrono scelto.
\end{narratore}

%\begin{narratore}
%In termini di Gioco di Ruolo canonico i Patroni sono divinità minori, creature di altri piani estremamente potenti e come tali \emph{affrontabili}, vedi ad esempio scheda di Tàhil. I Patroni della Genesi invece potendo creare i Patroni sono Dei Maggiori.
%\end{narratore}\end{changemargin}

\medskip

\textbf{Tabella Energia - Elementi}\index[Tabelle]{Tabella degli Elementi}

\medskip

\noindent\begin{tabularx}{\linewidth}{lXX}
	\toprule
\rowcolor{gray!20}\textbf{Positivi} (+1) & \textbf{Neutrali} (0) & \textbf{Negativi} (-1)\\
\toprule
En. Positiva& Fuoco& En. Negativa\\
\rowcolor{gray!20}Luce & Freddo& Vuoto\\
& Suono & \\
\rowcolor{gray!20}& Elettricità &
\end{tabularx}

\medskip

Nessun Patrono è completamente buono o completamente malvagio. Come la natura di un uomo è sfaccettata anche a seconda delle occasioni, così un Patrono potrà avere anomali slanci positivi o negativi. Ricordati che i Patroni sono Dei ma generati dalle emozioni e desideri.

\subsubsection{Miracoli, Interventi e Prodigi}\index{Miracoli, Interventi e Prodigi}

In un mondo dove le divinità sono così capricciose, volubili e assetati di devoti fa loro gioco dimostrarsi generosi con coloro che possono poi diffondere i loro Tratti.

Un favore chiesto ad un Patrono ha sempre un prezzo non ovvio ne scontato. Il Narratore deve valutare attentamente la supplica del personaggio e giudicare se la richiesta è pertinente con i Tratti del Patrono, in caso positivo tirare 1d100 e fare meno della metà del punteggio più alto di Tratto in comune con il Patrono. Oppure decidere autonomamente secondo il corso dell'avventura.

\addcontentsline{toc}{subsection}{Elenco Patroni}

%\subsubsection{Elenco Patroni}\label{elencopatroni}\index{Elenco Patroni}

\subsubsection{Ljust}\label{ljust}\index{Empatia}\hypertarget{ljust}{}

\index{Ljust}

\begin{enfasi}{
Solo la luce che uno accende a se stesso, risplende in seguito anche per gli altri. (Arthur Schopenhauer)
}\end{enfasi}


La Dama della Luce, colei che irradia calore e amore. Generatrice delle pulsioni d'amore, protezione, gentilezza, gioia e perdono. Racchiude in sé l'aspetto protettivo di una madre, la forza e l'audacia di una combattente, la passionalità di una giovane amante, l'allegria e la ricerca del nuovo, la fantasia di una bambina. Ljust incarna la bellezza della vita ed ogni creatura che la contempla vede quella che per lei è la massima armonia e cade prona al suo fascino.

Ljust può essere scelta solo da un personaggio con 4 Tratti in comune con lei, fondamentalmente si nasce per essere Devoti di Ljust. Nel corso delle ere Ljust decise di selezionare, scegliere e premiare le creature che più mostravano in modo innato e profondo amore per la vita, curiosità per il nuovo, forza incrollabile, dedizione, fiducia, rispetto e cura degli altri donando loro i poteri e la possibilità di studiare e crescere come Allieve della Luce. Queste Allieve devono seguire la regola degli 8 Passi.

\begin{tikzpicture}[remember picture, overlay]
	\node[anchor=center, opacity=0.3] at (0.45\columnwidth,0)
	{\includegraphics[width=1\columnwidth]{immagini/patroni-stella_luminosa.pdf}};
\end{tikzpicture}

\begin{itemize}[leftmargin=*] \setlength{\itemsep}{0pt}
\item \textbf{Simbolo}: Una stella a 8 punte con 8 raggi luminosi
\item \textbf{Caratteristica}(Devoto): Saggezza o Carisma
\item \textbf{Tratti}: Compassionevole, Testardo, Coraggioso, Estroverso, Altruista, Leale, Paziente. Il Devoto di Ljust ha 4 Tratti in comune con il Patrono.
\item \textbf{Manifestazione}: luce dorata inonda l'incantatore.
\item \textbf{Somma dei Tratti in comune a 5 punti}: puoi lanciare l'incantesimo Luce come Reazione, 3 volte al giorno
\item \textbf{Somma dei Tratti in comune a 10 punti}: guadagni un +2 ai Tiri Salvezza su Tempra

\item \textbf{Somma dei Tratti in comune a 15 punti}: una armatura di luce ti protegge, guadagni un +2 a tutti i Tiri Salvezza e Difesa, l'effetto è permanente.
\item \textbf{Somma dei Tratti in comune a 20 punti}: puoi lanciare l'incantesimo Bagliore Solare. 1 volta al giorno.
\item \textbf{Energia/B} (Seguace/Devoto): Energia Positiva, Luce
\item \textbf{Vantaggio} (Devoto): Ogni volta che fai una Cura magica curi un Punto Ferita in più.
\item \textbf{Liste Magia Privilegiate}(Seguace/Devoto): Cura, Abiurazione
\item \textbf{Arma Preferita}: Spada Bastarda
\item \textbf{Regola}: Accettare l'invito ad un ballo
\end{itemize}

\medskip

\textbf{Gli 8 Passi delle Allieve}\index{8 Passi delle Allieve}\index{Allieve}

\medskip

Le Allieve della Luce sono una gruppo segreto di Devote che per totale affinità con Ljust hanno intrapreso il difficile percorso del bene e dell'amore. E' tra i gruppi più antichi fondati sulla Terra. Le Allieve, 99 come numero massimo, ma purtroppo spesso meno numerose, sono Devote di Ljust e devono seguire gli 8 Passi della Luce

\begin{enumerate}[leftmargin=*] \setlength{\itemsep}{0pt}
\item Ama e proteggi con tutta te stessa, con totale e sincera dedizione chi hai attorno a te.
\item Non lasciare che la tua inazione generi sofferenza.
\item Si un punto di paragone. Fai che la tua Luce elevi le persone che hai intorno e possano vedere in te speranza, serenità, calma, protezione e sicurezza.
\item Usa l'intelligenza, la furbizia e l'arguzia. Si lungimirante e risoluta nell'azione.
\item La tua opera è per il bene comune. Fa che la tua Luce sia sempre alta ed intensa.
\item Non cercare altra Luce se non la tua e quella delle tue sorelle.
\item Sii luminosa ma non accecare chi è intorno a te.
\item Sii la differenza tra la disperazione e la speranza.
\end{enumerate}

Le Allieve hanno costruito un ballo armonioso trasformando in danza i passi della loro Regola.

Esistono anche Allieve di altro genere, rari ma storicamente accertati.

\subsubsection{Calicante}\index{Calicante}\label{calicante}\index{Egoismo}\hypertarget{calicante}{}

\begin{enfasi}{
La superstizione è la religione degli spiriti deboli. (Edmund Burke)
}\end{enfasi}

E' oscuro, gelido e arrabbiato. Racchiude in sé odio, violenza, distruzione, vendetta e perenne insoddisfazione. Raccoglie la personalità capricciosa e scontenta di un bambino, la noia violenta e sadica di un giovane uomo, la forza distruttiva di un uragano e la rabbia di un combattente che non ha più nulla da perdere. Calicante solo con la presenza mette a disagio, ti fa sentire in pericolo, affascina ma con le armi della paura e della incostanza.

Calicante può essere scelto solo dai personaggi che hanno 4 Tratti in comune con lui. I suoi Devoti sono i migliori assassini, sua professione più affine. Coloro che mostrano il maggiore sprezzo del pericolo e della vita altrui. I suoi prediletti sono coloro che sono temuti, odiati, coloro che sono violenti e crudeli ma mortalmente efficienti e decisivi in ogni situazione di combattimento.

\begin{itemize}[leftmargin=*] \setlength{\itemsep}{0pt}
\item \textbf{Simbolo}: Un turbine nero
\item \textbf{Caratteristica}: Forza o Destrezza
\item \textbf{Tratti}: Ambizioso, Disonesto, Vendicativo, Cinico, Dissoluto, Arrogante, Avaro. Il Devoto di Calicante ha 4 Tratti in comune con il Patrono
\item \textbf{Manifestazione}: spada grondante di sangue nero
\item \textbf{Somma dei Tratti in comune a 5 punti}: Puoi lanciare l'incantesimo Oscurità. Una volta al giorno
\item \textbf{Somma dei Tratti in comune a 10 punti}: La tua arma si ammanta di ombra. Guadagni un +2 al Tiro per Colpire e +1d4 di danno da Vuoto per 2d6 round, Una volta al giorno.
\item \textbf{Somma dei Tratti in comune a 15 punti}: Crei 4 dardi di Vuoto, ogni dardo fa 2d6 di danno, colpisce automaticamente entro 18 metri. Una volta al giorno.
\item \textbf{Somma dei Tratti in comune a 20 punti}: Crei una zona di energia protettiva attorno a te nel raggio di 3 metri, dimezzi tutto il danno che ricevi, non è possibile recuperare Punti Ferita nell'area. Durata 10 minuti consecutivi, 1 volta al giorno.
\item \textbf{Energia/M}: Energia Negativa, Vuoto
\item \textbf{Vantaggio}: +2 TS contro Lista Incantamento
\item \textbf{Liste Magia Privilegiate}: Fuoco, Necromanzia
\item \textbf{Arma Preferita}: Machete
\item \textbf{Regola}: Mai lasciare impunita una offesa diretta
\end{itemize}

\subsubsection{Atmos}\index{Atmos}\label{atmos}\index{Indifferenza}\hypertarget{atmos}{}

\begin{enfasi}{
Che cos'è dunque il tempo? Se nessuno me lo chiede, lo so; se voglio spiegarlo a chi me lo chiede, non lo so più. (Agostino da Ippona)
}\end{enfasi}

Il custode del Tempo e della Torre dell'Orologio, come ha avviato il tempo e la creazione dei nuovi Patroni così fermerà la sfida fra loro ed i Patroni sopravvissuti saranno giudicati, le loro opere valutate e Ljust o Calicante ne trarranno giovamento. Come una sfida da una singola moneta di rame nuovi Patroni, nuovi ideali saranno creati e noi, piccole creature vedremo nascere nuove civiltà e regni fiorenti. La storia è poco nota, solo i pochi Devoti di Atmos, scribi e studiosi della biblioteca del Tempo, conoscono il segreto e lo scorrere del tempo e della gara, gli altri, ignoranti, vivranno il loro tempo con un padrone sicuramente guidato da un Patrono.

Atmos, il Patrono del Tempo è il custode della storia, è colui che tiene traccia degli infiniti e più mondi che sono stati creati.

Atmos ha il potere unico e riservato solo a lui di poter bandire dal creato un Patrono qualora questo diventasse troppo forte e minacciasse Calicante e Ljust. Atmos ha già usato questo potere in passato. Atmos sia per la sua natura totalmente neutrale sia per il suo ruolo non si è mai schierato.

Tutti i Patroni temono Atmos per il suo potere, il più terribile per loro, ovvero il loro alienamento, l'oblio, la dimenticanza, l'essere distolti dal tempo e dalla sfida.

Per essere un Devoto di Atmos al momento del rito è necessario che il futuro Devoto possieda almeno quattro Tratti in comune con lui, amare la storia e la conoscenza.

Vestito di un morbido saio marrone e calzari di cuoio si muove tra gli infiniti scaffali della Biblioteca del Sapere con sempre uno strano misuratore del tempo appeso alla vita.


\begin{itemize}[leftmargin=*] \setlength{\itemsep}{0pt}
\item \textbf{Simbolo}: Un libro bianco con un orologio da taschino appoggiato sopra
\item \textbf{Caratteristica}: Intelligenza o Saggezza
\item \textbf{Tratti}: Indeciso, Prudente, Intransigente, Paziente, Vendicativo, Curioso, Avaro. Il Devoto di Atmos ha 4 Tratti in comune con il Patrono.
\item \textbf{Manifestazione}: l'incantesimo si sviluppa come a rallentatore, è solo un effetto illusorio
\item \textbf{Somma dei Tratti in comune a 5 punti}: Conosci sempre la data esatta e l'ora.
\item \textbf{Somma dei Tratti in comune a 10 punti}: Hai una intuizione innata per la conoscenza. Hai +1d6 alle prove di Conoscenza
\item \textbf{Somma dei Tratti in comune a 15 punti}: Puoi lanciare l'incantesimo Globo di Invulnerabilità, 1 volta al giorno.
\item \textbf{Somma dei Tratti in comune a 20 punti}: Ogni qual volta che devi fare una prova di Arcana puoi prendere il 18 come prendessi 10
\item \textbf{Energia/N}: Suono, Freddo
\item \textbf{Vantaggio}: Sai sempre che ora è
\item \textbf{Liste Magia Privilegiate}: Divinazione, Abiurazione
\item \textbf{Arma Preferita}: Mazza leggera
\item \textbf{Regola}: Non accettare o dare un compenso se non meritato
\end{itemize}

\subsubsection{Lynx}\index{Lynx}\label{lynx}\index{Responsabilità}\hypertarget{lynx}{}

\begin{enfasi}{
Le persone non fanno i viaggi, sono i viaggi che fanno le persone. (John Steinbeck)}
\end{enfasi}

Patrono dei Portali, è sceglibile solo da personaggi che abbiano almeno 3 Tratti in comune. E' il primo Patrono generato da Ljust e Calicante, creato per proteggere la Terra dagli attacchi esterni.

Serio, occhi gelidi di un azzurro chiarissimo è il Custode dei Portali e di ciò che è Oltre. Letale guardiano per chi cerca di passarli senza permesso, guida attenta per chi chiede il suo aiuto ed il suo permesso. Si fa scudo delle sue cicatrici per allontanare tutti. E' il solitario controllore del mondo.

I suoi Devoti sono i viaggiatori per eccellenza, coloro che presidiano e proteggono la Terra da ciò che è alieno, da ciò che potrebbe disturbare la creazione.

\begin{itemize}[leftmargin=*] \setlength{\itemsep}{0pt}
\item \textbf{Simbolo}: Un portale sull'oscurità
\item \textbf{Caratteristica}: Destrezza o Intelligenza
\item \textbf{Tratti}: Testardo, Coraggioso, Cinico, Intransigente, Vendicativo, Estroverso, Vanitoso
\item \textbf{Manifestazione}: come se il panorama non avesse più orizzonte
\item \textbf{Somma dei Tratti in comune a 5 punti} punti: Una volta al giorno puoi eseguire una Azione di Movimento in più
\item \textbf{Somma dei Tratti in comune a 10 punti}: Puoi lanciare Porta Dimensionale una volta al giorno
\item \textbf{Somma dei Tratti in comune a 15 punti}: Puoi lanciare l'incantesimo Esilio, 1 volta al giorno, DC 30.
\item \textbf{Somma dei Tratti in comune a 20 punti}: Puoi teletrasportarti per 500km al giorno (anche più teletrasporti o teletrasportati purché la somma totale non superi 500km)
\item \textbf{Energia/N}: Fuoco, Elettricità
\item \textbf{Vantaggio}: Sei considerato di una taglia maggiore quando devono spostarti o spingerti
\item \textbf{Liste Magia Privilegiate}: Evocazione, Acqua
\item \textbf{Arma Preferita}: Spada Corta
\item \textbf{Regola}: Non lasciare un ambiente non esplorato
\end{itemize}

\subsubsection{Gradh}\index{Gradh}\label{gradh}\index{Coraggio}\hypertarget{gradh}{}

\begin{enfasi}
L'uomo che ha cessato di avere paura ha cessato di preoccuparsi. (Francis Herbert Bradley)
\end{enfasi}

Il primo Patrono creato da Atmos sotto la guida di Ljust e l'influenza di Calicante.

Gradh racchiude in sé l'istinto innato alla protezione, alla difesa ed alla cura propri di Ljust. Gradh è quanto di più simile e profondamente legato a Ljust sia stato generato. Lui è equilibrio, razionalità ed empatia.
Dove vi è difesa, cura e protezione vi è Gradh.

Ma Calicante non poteva permettere la creazione di un Patrono totalmente votato ad Ljust e così infuse in Gradh la freddezza della vendetta e la furia della rabbia. Ecco che allora Gradh nell'atto di difendere l'umanità, spesso la deve in primis proteggere da sé stesso.

Gradh non ama sfidare apertamente Cattalm perché sa che farebbe esattamente il suo gioco, ecco che con astuzia cerca di attirarlo nel suo terreno di gioco, dove nessuna vita sarà in pericolo e lì da sfoggio a della sua superiorità strategica e di combattimento.

Passionale e freddo è forse il Patrono più umano del pantheon attuale. Il suo sguardo carismatico e protettivo può divenire freddo e tagliente quando è preda della furia della battaglia o della vendetta. Gradh ama studiare il mondo attorno a sé e passare inosservato. Spesso si nasconde fra la gente e \emph{vive} la sua vita umana, ma non si lascia avvicinare veramente da nessuno. Gradh attira a sé con la stessa facilità con cui lui allontana da se.

Il Devoto di Gradh è fiero ed orgoglioso, indomito e protettivo, ed addolorato, perché per quanto si sforzi di portare equilibrio e pace il male continua sempre a prosperare.

\begin{itemize}[leftmargin=*] \setlength{\itemsep}{0pt}
\item \textbf{Simbolo}: Uno scudo con incise sopra due spirali intrecciate.
\item \textbf{Caratteristica}: Forza
\item \textbf{Tratti}: Coraggioso, Vanitoso, Arrogante, Gentile, Invidioso, Leale, Sospettoso
\item \textbf{Manifestazione}: due spire una nera come ombra ed una lucente come scintilla circondano la tua arma intrecciandosi
\item \textbf{Somma dei Tratti in comune a 5 punti} punti: Puoi lanciare l'incantesimo Cura Ferite da 3 Punti Magia, ma ti causa 1d6 di danno. 1 volta al giorno
\item \textbf{Somma dei Tratti in comune a 10 punti}: Per 10 minuti consecutivi hai un bonus di +1d6 Tiro Salvezza su Riflessi e Tempra. Una volta al giorno.
\item \textbf{Somma dei Tratti in comune a 15 punti}: Emani un aura che concede a tutti i tuoi compagni entro raggio 3 metri un +2 Tiri Salvezza. Una volta al giorno, per 30 minuti consecutivi
\item \textbf{Somma dei Tratti in comune a 20 punti}: Lanci l'incantesimo Palla di Fuoco. L'incantesimo causa 60 di danno da energia negativa. DC 25 Riflessi per dimezzare. 2 volte al giorno
\item \textbf{Energia/N}: Energia Positiva - Energia Negativa
\item \textbf{Vantaggio}: +2 Consapevolezza
\item \textbf{Liste Magia Privilegiate}: Abiurazione, Invocazione
\item \textbf{Arma Preferita}: Mazza flangiata
\item \textbf{Regola}: Non permettere ad un immondo di camminare sulla Terra
\end{itemize}

\subsubsection{Atherim}\index{Atherim}\label{atherim}\index{Onestà}\hypertarget{atherim}{}

\begin{enfasi}{
Colui al quale confidate il vostro segreto, diventa il padrone della vostra libertà. (François de La Rochefoucauld)

\medskip

Non c'è nulla di nascosto che non sarà svelato, né di segreto che non sarà conosciuto. (Luca, 12, 1-7)
}\end{enfasi}

Il Patrono custode. Molti vedono nel seno generoso di Atherim un segno di voluttà e passione. Si lasciano incantare dalla sua procace bellezza e non vedono gli occhi di cristallo che incutono timore a chi osa anche solo pensare di avvicinarla.

Atherim è la custode dei sogni e delle speranze, colei alla quale affidare, come ad una madre, i desideri. E' il Patrono dei Bambini, dei Segreti e delle Levatrici.

Dal sorriso allegro e dall'animo buono sarà sempre pronta ad aiutarti a realizzare i tuoi sogni. E come una madre Atherim protegge e custodisce i segreti e le passioni. Atherim è muta. E' colei che custodisce per sempre, dentro il suo animo i segreti.

Il Devoto di Atherim si prende a cuore coloro che hanno fatto una promessa, punisce chi le infrange e chi svela i segreti. Molti Devoti di Atherim sono diplomatici, notai e levatrici.

\begin{itemize}[leftmargin=*] \setlength{\itemsep}{0pt}
\item \textbf{Simbolo}: Una mano di donna guantata che tiene un'ampolla ricca di flussi
\item \textbf{Caratteristica}: Saggezza
\item \textbf{Tratti}: Sospettoso, Compassionevole, Altruista, Intransigente, Coraggioso, Entusiasta, Vanitoso
\item \textbf{Manifestazione}: un silenzio sereno e tranquillizzante cala attorno all'incantatore
\item \textbf{Somma dei Tratti in comune a 5 punti}: Puoi aggiungere 1d6 ad un Tiro salvezza dopo averlo tirato ma prima di sapere se ha avuto successo o meno. Una volta al giorno, come Reazione.
\item \textbf{Somma dei Tratti in comune a 10 punti}: Guadagni 30 Punti Ferita temporanei. Durata 1 ora, una volta al giorno, come Azione Immediata.
\item \textbf{Somma dei Tratti in comune a 15 punti}: Puoi lanciare l'incantesimo \hyperlink{Zona di Verità}{Zona di Verità} 3 volte al giorno, senza Tiro Salvezza.
\item \textbf{Somma dei Tratti in comune a 20 punti}: Ogni pozione che bevi ha il doppio come durata o effetto se immediata.
\item \textbf{Energia/B}: Energia Positiva, Elettricità
\item \textbf{Vantaggio}: Riduci il Sanguinamento di 1 a fine round
\item \textbf{Liste Magia Privilegiate}: Ammaliamento
\item \textbf{Arma Preferita}: Pugnale
\item \textbf{Regola}: Non rivelare un segreto confidato
\end{itemize}

\subsubsection{Belevon}\index{Belevon}\label{belevon}\index{Invidia}\hypertarget{belevon}{}

\begin{enfasi}{
Dove c'è un uomo, c'è anche la bugia. (Robert Louis Stevenson)

\medskip

Nessuno ha una memoria tanto buona da poter essere un perfetto bugiardo. (Abraham Lincoln)
}\end{enfasi}

E' il Patrono che meglio incarna la bugia e la finzione al fine di un proprio tornaconto. Lui ama solo se stesso. E' un narcisista che si circonda solo di persone che lo assecondano e lo adulano. Aborrisce la solitudine ma allo stesso tempo odia essere toccato da qualcuno.

E' sempre alla ricerca di nuove cose, di oggetti meravigliosi che scambia e ricambia con altri oggetti. Gli piace discutere e mercanteggiare fino ad ottenere sempre quello che vuole a costo della vita altrui.

Belevon è un orrendo e deforme Patrono adorato da abiette creature delle caverne più profonde, invidiose di quanto posseduto dagli altri.

Il Devoto di Belevon è ben descritto da un uomo lucertola circondato da chincaglieria e resti umani.

\begin{itemize}[leftmargin=*] \setlength{\itemsep}{0pt}
\item \textbf{Simbolo}: Una gabbia dorata
\item \textbf{Caratteristica}: Intelligenza
\item \textbf{Tratti}: Invidioso, Ambizioso, Dissoluto, Disonesto, Compassionevole, Paziente, Altruista
\item \textbf{Manifestazione}: come se le sbarre dorate di una gabbia si intrecciassero attorno all'incantatore
\item \textbf{Somma dei Tratti in comune a 5 punti} punti: Puoi lanciare l'incantesimo Prestidigitazione, 3 volte al giorno.
\item \textbf{Somma dei Tratti in comune a 10 punti}: Acquisisci la capacità lanciare l'incantesimo \hyperlink{Immagine Maggiore}{Immagine Maggiore} una volta al giorno.
\item \textbf{Somma dei Tratti in comune a 15 punti}: Puoi lanciare l'incantesimo Allucinazione Mortale. 1 volta al giorno
\item \textbf{Somma dei Tratti in comune a 20 punti}: Toccando un oggetto vieni a conoscenza per sommi capi della storia di chi l'ha creato. Una volta al giorno. Costa 3 Azioni.
\item \textbf{Energia/N}: Fuoco, Suono
\item \textbf{Vantaggio}: Vedi Abilità Fortunato
\item \textbf{Liste Magia Privilegiate}: Illusione
\item \textbf{Arma Preferita}: Picca leggera
\item \textbf{Regola}: Contratta sul prezzo che sia in acquisto o vendita
\end{itemize}

\subsubsection{Cattalm}\index{Cattalm}\label{cattalm}\index{Cinismo}\hypertarget{cattalm}{}

\begin{enfasi}{
Non è l'essere arrabbiati che conta, è l'essere arrabbiati per le cose giuste. Le dissi: guardalo dalla prospettiva darwiniana. La rabbia serve a renderti efficiente. Questa è la sua funzione per la sopravvivenza. Ecco perché ti è stata data. Se ti rende inefficiente, mollala come una patata bollente. (Philip Roth)
}\end{enfasi}

Generato direttamente da Calicante, come risposta alla creazione di Gradh da parte di Ljust, è pura distruzione, caos ed entropia. Cattalm si prefigura il solo scopo di distruggere, portare caos e malattie, terremoti ed alluvioni.

Cattalm è tra i pochi Patroni che osa sfidare apertamente Gradh e lo fa con gioia perché sa che la loro battaglia altro non farà che portare ulteriore distruzione. Cattalm accetta ed invita ad essere suo Devoto ogni creatura capace di odio, capace di distruggere e ferire. Molti suoi Devoti sono creature mostruose o aberrazioni.

Cattalm invece è tra i Patroni più meravigliosi, con una candida pelle lucente, ali di piuma soffice ed una leggera armatura argentata. Per quanto i lineamenti delicati ne facciano un essere bellissimo per quanto ambisca alla distruzione.

Cattalm adora il caos che manifesta nei modi più violenti con terremoti, alluvioni, maremoti, malattie se non direttamente piogge infuocate. Non agisce quasi mai direttamente ma lascia che caos e distruzione lavorino per lui.

Ljust non poteva non intervenire nella creazione di un Patrono così esplicitamente malvagio e, di nascosto da Calicante, instillò in Cattalm l'amore e protezione per i bambini. Cattalm distrugge, avvelena, indebolisce ma non i bambini, neanche indirettamente, piuttosto si attiva lui stesso per annullare i malefici causati dalla sua natura.

Ogni qual volta succede una calamità si suole dire che \emph{Cattalm ha battuto il piede}.

\begin{itemize}[leftmargin=*] \setlength{\itemsep}{0pt}
\item \textbf{Simbolo}: Un'onda gigante che sovrasta la costa
\item \textbf{Caratteristica}: Forza
\item \textbf{Tratti}: Cinico, Arrogante, Ambizioso, Intransigente, Dissoluto, Sospettoso, Paziente
\item \textbf{Manifestazione}: il rombo del tuono
\item \textbf{Somma dei Tratti in comune a 5 punti} punti: Attraverso le tue armi indebolisci l'avversario designato. A seguito di un colpo critico puoi aumentare di un livello l'affaticamento. Una volta al giorno come Reazione.
\item \textbf{Somma dei Tratti in comune a 10 punti}: Il tuo tocco imputridisce cibo (fino a 50kg/Ingombro 10) e acqua (un cubo con uno spigolo di 10 m). Una volta al giorno
\item \textbf{Somma dei Tratti in comune a 15 punti}: Il tuo sguardo accieca di collera. Lanci l'incantesimo Confusione, ma l'unico risultato possibile è che i target attaccano dei soggetti a caso. DC 25. 1 volta al giorno
\item \textbf{Somma dei Tratti in comune a 20 punti}: Lanci l'incantesimo Cono di Freddo 60 di danno, ma il danno è da Vuoto. DC 25. Una volta al giorno
\item \textbf{Energia/M}: Energia Negativa, Vuoto
\item \textbf{Vantaggio}: aumenta di 10 i PF necessari ad ucciderti
\item \textbf{Liste Magia Privilegiate}: Fuoco
\item \textbf{Arma Preferita}: Ascia da battaglia
\item \textbf{Regola}: Non fare buone azioni senza un tornaconto
\end{itemize}

\subsubsection{Efrem}\index{Efrem}\label{efrem}\index{Rispetto}\hypertarget{efrem}{}

\begin{enfasi}{
Non deviare dalla natura ed il formarci sulle sue leggi e sui suoi esempi, è sapienza. (Lucio Anneo Seneca)
}\end{enfasi}

E' il Patrono di chi fa della natura la propria casa. Incarna in sé gli aspetti più puri della natura stessa, aggressivo come solo i felini più letali sanno essere; ma anche selvaggio come le radure più nascoste e rigorosa come solo la natura può essere.

Efrem si prefigge di difendere la Natura dalla contaminazione dell'uomo, da questa specie infestante che distrugge tutto ciò che incontra.

I Devoti di Efrem. chiamati anche druidi\index{Druido}, sono legati maggiormente all'elemento naturale. Manipolano la magia principalmente elementale e si difendono o attaccano usando anche animali e creature naturali. Si dice che i più potenti costringano anche i Draghi alla ubbidienza.

I Devoti di Efrem hanno l'obiettivo supremo di proteggere gli animali e le piante, i luoghi e tutto ciò che è naturale e non artificiale. Solitamente solitario e scontroso non riesce a capire il perché dell'odio che, dal suo punto di vista, l'uomo scarica sulla Terra.

Un Devoto di Efrem rispetta la vita come la morte, nel processo naturale che è l'evoluzione ed il ciclo vitale. A volte decide di stabilirsi in un certo ambiente e lo elegge come suo territorio e come fosse la sua casa lo protegge. Altre volte decide di essere ramingo ed intervenire in tutto il mondo per proteggere le sue amate piante ed animali.

Nelle terre più desolate, nelle regioni più naturali i Devoti di Efrem costruiscono utopie tra umanoidi e animali, dove l'equilibrio viene mantenuto con il sangue di qualsiasi che si ribelli al loro volere.

\begin{itemize}[leftmargin=*] \setlength{\itemsep}{0pt}
\item \textbf{Simbolo}: Una staffa con un rampicante attorcigliato attorno
\item \textbf{Caratteristica}: Costituzione
\item \textbf{Tratti}: Leale, Indeciso, Prudente, Impulsivo, Testardo, Paziente, Ambizioso
\item \textbf{Manifestazione}: spire di foglie avvolgono l'arma
\item \textbf{Somma dei Tratti in comune a 5 punti} punti: Il tuo tocco rende docili gli animali non magici. Tiro Salvezza su Volontà 20 per resistere. 3 volte al giorno. Costo 2 Azioni.
\item \textbf{Somma dei Tratti in comune a 10 punti}: Guadagni un +1d6 a tutte le prove di Sopravvivenza che si effettuano in un ambiente naturale.
\item \textbf{Somma dei Tratti in comune a 15 punti}: Puoi lanciare l'incantesimo Bacche Benefiche 1 volta al giorno. Ogni bacca cura 1d6 Punti Ferita e rimuove le malattie o veleni non magici.
\item \textbf{Somma dei Tratti in comune a 20 punti}: Il tuo tocco è quello del padrone. Puoi ammansire creature anche magiche come Aberrazioni o Draghi che tocchi. Tiro Salvezza su Volontà DC 30. Una volta al giorno. Costo 2 Azioni
\item \textbf{Energia/N}: Elettricità, Suono
\item \textbf{Vantaggio}: +1d6 a Gestire Animali
\item \textbf{Liste Magia Privilegiate}: Animali e Piante ed una Liste Magia Elementale.
\item \textbf{Arma Preferita}: Bastone
\item \textbf{Regola}: La Natura è sempre la tua prima scelta
\end{itemize}

\subsubsection{Erondil}\index{Erondil}\label{erondil}\hypertarget{erondil}{}\index{Superbia}

\begin{enfasi}{
I buoni ragionamenti sono più forti di due mani robuste. (Sofocle)
}\end{enfasi}

Patrono di Terra e Aria, Erondil è il Signore degli elementi più concreti e razionali. Colui che dotato di infinito potere e razionalità concede ai suoi Devoti il potere della manipolazione della terra, il dono di creare dal semplice fango costruzioni gigantesche e di millenaria forza. Conclude le sue opere con attenzione e precisione.
Pur con fatica perché se il risultato finale non lo soddisfa scatena i suoi fulmini per distruggerlo all'istante. Perfezionista ed incontentabile, difficilmente qualcosa è esattamente come lui se la immaginava.

Ordinato ed esuberante è il signore delle tempeste, dei tuoni e dei fulmini, dei terremoti e distruzioni. Ama circondarsi del fragore del tuono, del rombo della terra che si sgretola. Sa essere distruttivo verso coloro che non rispettano la Terra.
Ha braccia e petto ricoperti da tatuaggi quasi argentei che narrano le leggende della Terra e dell'Aria.

I Devoti di Erondil sono gli ingegneri dell'impossibile.

\begin{itemize}[leftmargin=*] \setlength{\itemsep}{0pt}
\item \textbf{Simbolo}: un castello di sabbia con un fulmine sopra
\item \textbf{Caratteristica}: Saggezza
\item \textbf{Tratti}: Arrogante, Vendicativo, Ambizioso, Compassionevole, Entusiasta, Leale, Avaro
\item \textbf{Manifestazione}: suono di tempesta e rombo di frana
\item \textbf{Somma dei Tratti in comune 5 punti}: Non temi più le cadute. Puoi lanciare l'incantesimo Caduta Piuma 3 volte al giorno, solo su di te.
\item \textbf{Somma dei Tratti in comune a 10 punti}: Il tuo tocco plasma la pietra. Puoi lanciare l'incantesimo Passa Porta 1 volta al giorno.
\item \textbf{Somma dei Tratti in comune a 15 punti}: Puoi scagliare l'incantesimo Fulmine dalle tue mani. Tiro Salvezza Riflessi DC 30 per dimezzare. Costo 2 Azioni.
\item \textbf{Somma dei Tratti in comune a 20 punti}: Sei in grado di creare una fossa profondissima (1km) sotto il tuo avversario (taglia fino a grande). Tiro Salvezza Riflessi 30 o cadere ogni round di 100 metri. Una volta al giorno. Dopo 1 minuto la fossa si chiude con chi c'è dentro. Costo 2 Azioni.
\item \textbf{Energia/N}: Suono, Elettricità
\item \textbf{Vantaggio}: Riduzione 10 ai danni da Elettricità
\item \textbf{Liste Magia Privilegiate}: Aria, Terra
\item \textbf{Arma Preferita}: Martello da guerra
\item \textbf{Regola}: Non devi consentire la distruzione di monumenti architettonici
\end{itemize}

\subsubsection{Gaya}\index{Gaya}\label{gaya}\hypertarget{gaya}{}\index{Generosità}

\begin{enfasi}{
{Splendore del giorno concluso, che mi sollevi e mi colmi,

ora profetica, ora che il passato riadduci!

E mi gonfi la gola, te, divino egualitarissimo,

voi, terra e vita finché brilli l'ultimo raggio, io canto. (Song at Sunset, Walt Whitman)}
}\end{enfasi}

Patrono di Acqua e Fuoco, nelle profondità della terra, dove acqua e lava si incontrano, Gaya si diverte a dipingere. Adora circondarsi dei flussi di fuoco e acqua quasi a creare una danza in mezzo a loro. Adora i suoni della natura, l'infrangersi delle onde sugli scogli, il cadere delle gocce di pioggia sull'acciottolato, il borbottare di un fuoco scoppiettante.

Dipinge mescolando il caldo ed il freddo. L'acqua cristallina ed impetuosa al fuoco intrigante ed ardente. Gelosa del bello e delle arti tiene tutte le sue opere al sicuro in un ordine quasi maniacale e protette. Da vera artista utilizza gli elementi per far risplendere le meraviglie della natura. Gaya è la pittrice di tramonti e delle tempeste.

I Devoti di Gaya sono artisti volubili e sopra le righe. Sono coloro che ricreano la magia dell'alba o del tramonto o del mare in tempesta nelle loro opere, sono coloro che mettono poesia e follia nella normalità.

Ma Gaya ha anche un lato molto più subdolo e violento, una vena di follia malvagia che adora portare distruzione con fiamme e acqua. Nelle profondità delle caverne creature affini all'acqua od al fuoco adorano Gaya e uccidono chiunque non sia d'accordo con loro.

\begin{narratore}[Gaia ed Erondil]
\textbf{Gaia} ed \textbf{Erondil} sono come le due facce della stessa medaglia e sovraintendono agli elementi, Gaia acqua e fuoco ed Erondil Aria e Terra; agiscono come espressione diretta dei Patroni maggiori, sono piccole manifestazione del loro immenso potere.
\end{narratore}

\begin{itemize}[leftmargin=*] \setlength{\itemsep}{0pt}
\item \textbf{Simbolo}: un pennello sul cielo
\item \textbf{Caratteristica}: Intelligenza
\item \textbf{Tratti}: Altruista, Gentile, Sospettoso, Cinico, Invidioso, Disonesto, Arrogante
\item \textbf{Manifestazione}: spire di fuoco e acqua avvolgono all'incantatore
\item \textbf{Somma dei Tratti in comune a 5 punti} punti: Puoi creare fino a 5 litri di acqua o 1 litro di liquore di buona qualità. Una volta al giorno. Costo 2 Azioni.
\item \textbf{Somma dei Tratti in comune a 10 punti}: Il tuo metabolismo non teme il freddo. Resisti al Danno magico da freddo e sei immune a quello naturale.
\item \textbf{Somma dei Tratti in comune a 15 punti}: Puoi respirare sott'acqua come respiri l'aria. Resisti al danno da fuoco non magico
\item \textbf{Somma dei Tratti in comune a 20 punti}: Generi una pioggia di fuoco. Lanci l'incantesimo Colpo Infuocato, DC 25 una volta al giorno. Resisti il danno da fuoco magico.
\item \textbf{Energia/N}: Freddo, Fuoco
\item \textbf{Vantaggio}: Riduzione 10 ai danni da Fuoco
\item \textbf{Liste Magia Privilegiate}: Acqua, Fuoco
\item \textbf{Arma Preferita}: Tridente
\item \textbf{Regola}: Non impedirti di ascoltare della buona musica
\end{itemize}

\subsubsection{Krondal}\index{Krondal}\label{krondal}\index{Intransigenza}\hypertarget{krondal}{}

\begin{enfasi}{
La libertà, Sancio, è uno dei più preziosi doni che i cieli abbiano mai dato agli uomini; né i tesori che racchiude la terra né che copre il mare sono da paragonare ad essa; per la libertà, come per l'onore, si può e si deve mettere a repentaglio la vita. (Miguel de Cervantes)
}\end{enfasi}

Krondal il folle, Krondal l'assassino, Krondal il salvatore.

Questi e molti altri sono gli appellativi di Krondal, il Patrono che non potrai mai veramente capire completamente.

Krondal abbraccia lo spirito anarchico e libero nella maniera più assoluta. Tutti secondo Krondal devono fare solo ciò che vogliono.

Il suo motto è \emph{Nessuno conosce Nessuno} perché non puoi sapere il futuro e quello che ti aspetta.

Krondal nutre un profondo rispetto per la libertà e non può criticare le scelte, estreme o meno che vengano fatte, eppure per suo dettato divino perseguita intransigente per portare giustizia.

Un Devoto di Krondal è tipicamente una guardia del corpo, un protettore, lo sceriffo al quale non interessano i motivi della scelta ma che sa giudicare le azioni compiute.

\begin{itemize}[leftmargin=*] \setlength{\itemsep}{0pt}
\item \textbf{Simbolo}: Una spada tenuta verticalmente davanti a se
\item \textbf{Caratteristica}: Carisma
\item \textbf{Tratti}: Intransigente, Vanitoso, Arrogante, Sospettoso, Paziente, Ambizioso, Testardo
\item \textbf{Manifestazione}: il mantello o veste del Devoto diventa pulito e lucente
\item \textbf{Somma dei Tratti in comune a 5 punti} punti: Maledici il tuo avversario. Lanci una volta al giorno l'incantesimo \hyperlink{Scagliare Maledizione}{Scagliare Maledizione}. DC 20 per resistere.
\item \textbf{Somma dei Tratti in comune a 10 punti}: Non vuoi essere legato o ammanettato. Due volte al giorno puoi lanciare solo su te stesso Libertà di Movimento.
\item \textbf{Somma dei Tratti in comune a 15 punti}: La tua presenza toglie la vista agli avversari. Designa fino a 6 creature entro 9 metri, queste devono fare un Tiro Salvezza Tempra a DC 30 o essere ciechi solo nei tuoi confronti per 1d4 round.
\item \textbf{Somma dei Tratti in comune a 20 punti}: La tua arma è più efficace contro i nemici. Ogni creatura colpita deve fare un Tiro Salvezza Volontà DC 20 o rimanere paralizzata per 3 round. Una volta che la creatura riesce nel Tiro Salvezza non può essere più influenzata per le successive 24 ore. Una volta la giorno, attivare l'abilità costa 1 Azione e dura 1 minuto.
\item \textbf{Energia/B}: Energia Positiva, Fuoco
\item \textbf{Vantaggio}: +2 TS su Tempra
\item \textbf{Liste Magia Privilegiate}: Abiurazione
\item \textbf{Arma Preferita}: Spada lunga
\item \textbf{Regola}: Non permettere soprusi
\end{itemize}

\subsubsection{Ledyal}\index{Ledyal}\index{Laydel}\label{ledyal}\label{laydel}\hypertarget{ledyal}{} \hypertarget{laydel}{}\index{Gratitudine}\index{Rancore}

\begin{enfasi}{
L'anima di un uomo è immortale e incorruttibile. (Platone)

\medskip

Non è il corpo che mi definisce (anonima creatura)
}\end{enfasi}

E' il Patrono senza un volto preciso, senza una voce se non un canto. Mutevole nel corpo e senza una definizione chiara del suo essere. Si manifesta con un lungo mantello color rosso fuoco dal tessuto fatto da mille farfalle. Il suo tocco è vita e pace, protegge chi necessita dei suoi favori indipendentemente dal fatto che li chieda o meno. Desidera un mondo senza sofferenza, con solo felicità ed armonia. Sospettoso/a e profondamente introverso non crede a coloro che gli danno ragione. Ha il cuore pieno di vita e di bontà ma non ha un corpo con cui amare.

Ledyal ha anche una sorella gemella, o forse un'altra personalità, o forse sono lo stesso Patrono, nessuno le ha mai viste insieme. La \emph{gemella} \textbf{Laydel}\index{Laydel} non tollera la sofferenza, disprezza chi causa dolore, uccide senza timore qualunque creatura abbia peccato contro un innocente, chiunque abbia causato sofferenza.

\begin{itemize}[leftmargin=*] \setlength{\itemsep}{0pt}
\item \textbf{Simbolo}: Una farfalla che gronda sangue mentre vola
\item \textbf{Caratteristica}: Saggezza (Ledyal) - Forza (Laydel)
\item \textbf{Tratti Ledyal}: Entusiasta, Compassionevole, Prudente, Gentile, Curioso, Codardo, Testardo
\item \textbf{Tratti Laydel}: Vendicativo, Paziente, Ambizioso, Intransigente, Invidioso, Cinico, Arrogante
\item \textbf{Manifestazione}: come se un mantello di farfalle avvolgesse il Devoto
\item \textbf{Somma dei Tratti in comune a 5 punti} punti: Il tuo tocco è vita/attacco. 3 volte al giorno puoi toccare una creatura vivente e curarla/causare 1d6 Punti Ferita. Costo 2 Azioni (comprende anche l'Azione di tocco)
\item \textbf{Somma dei Tratti in comune a 10 punti}: Il tuo tocco è pace. Puoi lanciare 2 volta al giorno l'incantesimo Santuario.
\item \textbf{Somma dei Tratti in comune a 15 punti}: La tua aura protegge i tuoi compagni. Entro raggio 6 metri i tuoi compagni hanno un +4 alla Difesa ed un +2 ai Tiri Salvezza. Durata 10 minuti consecutivi, una volta al giorno. Costo 2 Azioni.
\item \textbf{Somma dei Tratti in comune a 20 punti}: Irradi una sfera curativa intorno a te. Ogni creatura nel raggio di 6 metri viene curata di 60 Punti Ferita. Una volta al giorno. In caso di Laydel l'effetto è opposto. Costo 2 Azioni
\item \textbf{Energia/B}: Energia Positiva, Elettricità
\item \textbf{Vantaggio}: +1d6 alle prove di Pronto Soccorso (Ledyal) oppure sei +4 TS contro Paura (Laydel)
\item \textbf{Liste Magia Privilegiate}: Cura o Invocazione
\item \textbf{Arma Preferita}: Manganello/Catena chiodata
\item \textbf{Regola}: Non permettere violenze contro il genere di una creatura
\end{itemize}

\subsubsection{Nethergal}\index{Nethergal}\label{nethergal}\index{Sincerità}\hypertarget{nethergal}{}

\begin{enfasi}{
I sogni sono risposte a domande che non siamo ancora in grado di formulare. (X-Files)

\medskip

Un sogno non interpretato è come una lettera non letta. (Talmud)
}\end{enfasi}

\medskip

Il Patrono Messaggero. Sulla piuma di un'oca vola la lettera di Nethergal. Rapida, impetuosa, diretta, Nethergal è la messaggera, colei alla quale affidare pensieri e scritti. Sarcastica e logorroica curioserà sui tuoi scopi, ti chiederà informazioni sugli scritti affidatole con esplicita franchezza ed avrà sempre qualcosa da ridire sul messaggio da portare ma sarà anche altrettanto diretta e precisa nel consegnarlo.

Nethergal non è solo chiacchiere e pettegolezzi, qualsiasi testo venga scritto lei lo conosce, non esiste codice o segreto scritto che lei non conosca o non sappia decifrare.

Il Devoto di Nethergal è un fine linguista, un esperto di indovinelli e rebus, un Devoto che a differenza di Atmos non si limita a custodire gli scritti ma ne diffonde la conoscenza.

Un Devoto di Nethergal è un maestro, una professoressa di lingue di un Collegio, un dotto esperto di mille argomenti, probabilmente vanitoso e qualche volta arrogante se sfidato sulle sue materie.

Nethergal conosce la locazione di ogni documento ed è probabilmente la chiave per comprendere cosa successe alla Freten. Pochissimi sanno che Nethergal era uno dei Patroni della prima venuta, quella incaricata di distruggere le infrastrutture di comunicazione, colei che cifrò i contenuti degli archivi elettronici nel più grande ransoware della storia.

Nethergal ha anche un altro ruolo è il Patrono dei sogni e delle visioni, divide questo compito con Sixiser che invece domina gli incubi.

\begin{itemize}[leftmargin=*] \setlength{\itemsep}{0pt}
\item \textbf{Simbolo}: una piuma bianca cangiante
\item \textbf{Caratteristica}: Destrezza
\item \textbf{Tratti}: Estroverso, Curioseffettuareo, Testardo, Vanitoso, Vendicativo, Arrogante, Paziente
\item \textbf{Manifestazione}: cascata di piume, un oca in volo
\item \textbf{Somma dei Tratti in comune a 5 punti} punti: Puoi inviare un messaggio di massimo 144 caratteri ad un soggetto che puoi vedere entro 50 metri senza essere udito/visto. Una volta all'ora. Costo 1 Azione. Il soggetto deve comprendere la lingua usata.
\item \textbf{Somma dei Tratti in comune a 10 punti}: Mettendo la mano su un libro ne apprendi il contenuto come se lo avessi letto. Un libro a settimana. Perdi le conoscenze così acquisite dopo una settimana. Tempo 1 round. La lingua scritta del tomo deve essere nota.
\item \textbf{Somma dei Tratti in comune a 15 punti}: Puoi volare, come omonimo incantesimo, 1 ora al giorno. Costo 2 Azioni.
\item \textbf{Somma dei Tratti in comune a 20 punti}: Comprendi ogni scritto che non sia magico o codificato.
\item \textbf{Energia/N}: Elettricità, Suono
\item \textbf{Vantaggio}: sai sempre dove è il nord magnetico
\item \textbf{Liste Magia Privilegiate}: Trasmutazione, Aria
\item \textbf{Arma Preferita}: Balestra leggera
\item \textbf{Regola}: Non distruggere un libro o una lettera
\end{itemize}

\subsubsection{Nedraf}\index{Nedraf}\label{nedraf}\index{Perseveranza}\hypertarget{nedraf}{}

\begin{enfasi}{
Pensa veramente di lottare per qualcosa a parte la sua sopravvivenza? (Matrix Revolutions, Film)

\medskip

Non importa quanto stretto sia il passaggio,

Quanto piena di castighi la vita.

Io sono il padrone del mio destino:

Io sono il capitano della mia anima (Invictus, William Ernest Henley)
}\end{enfasi}

Il Patrono Sopravvissuto, il vecchio lupo mai stanco che ha attraversato e combattuto innumerevoli battaglie. La sua carne è ferita, il suo corpo ricoperto di cicatrici di guerra e lividi ma nulla lo farà crollare. Tenacia, passione, esperienza e tanta rabbia rendono Nedraf non solo un combattente eccellente in qualsiasi occasione ma un conoscitore dell'ambiente attorno a sé. Grazie al suo impeccabile allenamento sa sfruttare al meglio le risorse a disposizione. Sa spronare con passione gli uomini a suoi ordini.
Nedraf rappresenta colui che vorresti sempre accanto in ogni battaglia.

Molti capitani di ventura e ufficiali al comando sono Devoti di Nedraf. Il Devoto di Nedraf non si arrende, non rinuncia, non abbandona i compagni ma non per questo è avventato o irrazionale nelle scelte.

\begin{itemize}[leftmargin=*] \setlength{\itemsep}{0pt}
\item \textbf{Simbolo}: una mano forte, avvolta in una benda sporca di sangue che brandisce una spada
\item \textbf{Caratteristica}: Costituzione
\item \textbf{Tratti}: Paziente, Vanitoso, Coraggioso, Intransigente, Entusiasta, Arrogante, Cinico
\item \textbf{Manifestazione}: si spande nell'aria odore di sangue e metallo
\item \textbf{Somma dei Tratti in comune a 5 punti} punti: Puoi portare armature leggere senza penalità alla Prova di Magia
\item \textbf{Somma dei Tratti in comune a 10 punti}: Acquisisci un punto bonus su una Lista armi. Può essere nota o meno
\item \textbf{Somma dei Tratti in comune a 15 punti}: Puoi portare armature medie senza penalità alla Prova di Magia ed Destrezza
\item \textbf{Somma dei Tratti in comune a 20 punti}: Acquisisci un punto bonus su una Lista armi. Può essere nota o meno
\item \textbf{Energia/B}: Energia positiva, Suono
\item \textbf{Vantaggio}: recuperi il doppio di PF quando riposi
\item \textbf{Liste Magia Privilegiate}: Ammaliamento, Terra
\item \textbf{Arma Preferita}: Spadone a due mani
\item \textbf{Regola}: Non abbandonare i compagni
\end{itemize}

\subsubsection{Nihar}\index{Nihar}\label{nihar}\index{Umiltà}\hypertarget{nihar}{}

\begin{enfasi}{
Che cos'è un eroe? È un individuo dotato di un grande talento e straordinario coraggio, che sa scegliere il bene al posto del male, che sacrifica se stesso per salvare gli altri, ma soprattutto... che agisce quando ha tutto da perdere e nulla da guadagnare. (Lo chiamavano Jeeg Robot, film)
}\end{enfasi}

E' il Patrono degli eroi per caso. Ponderato e tranquillo è amante del buon vino e del gozzovigliare. E' colui che non sceglieresti mai come compagno d'armi a causa del suo aspetto \emph{comune} e del suo atteggiamento goliardico. Ma poi al momento di esserci, di combattere, di far la differenza ecco che con un colpo fortunato risolve la sfida.

Ha le sembianze di un piccolo uomo, dai vestiti sfarzosi e ricercati e dall'espressione guardinga ed allegra. Si protegge sempre e a qualunque costo, mostrando al mondo esattamente ciò che il mondo vuole vedere. Controlla attentamente la realtà attorno a sé e anche se è sempre più facile vederlo con un calice in mano, se non ci si lascia ingannare dalle apparenze si noterà come i suoi occhi non perdano mai di vista il pericolo, il problema. Sta attento, non si fida di nulla e di nessuno. Ha fatto dei suoi difetti i suoi punti di forza.

\begin{itemize}[leftmargin=*] \setlength{\itemsep}{0pt}
\item \textbf{Simbolo}: Una daga appoggiata vicino ad un calice di vino
\item \textbf{Caratteristica}: Intelligenza
\item \textbf{Tratti}: Curioso, Coraggioso, Compassionevole, Vanitoso, Invidioso, Avaro, Crudele
\item \textbf{Manifestazione}: il suono di un brindisi o lo stappare di una bottiglia
\item \textbf{Somma dei Tratti in comune a 5 punti} punti: Puoi trasformare l'acqua in vino. Un litro al giorno. Costo 2 Azioni. 2 volte al giorno.
\item \textbf{Somma dei Tratti in comune a 10 punti}: Una Azione Immediata, ottieni un bonus di +2d6 ad una prova di Competenza in quel round. 3 volta al giorno.
\item \textbf{Somma dei Tratti in comune a 15 punti}: La tua arma leggera causa sempre un danno critico quando colpisci. Il bonus è sempre attivo.
\item \textbf{Somma dei Tratti in comune a 20 punti}: I manicaretti che prepari sono buonissimi. Chiunque si sazi con una pietanza da te preparata recupera 2d6 Punti Ferita e viene curato dai veleni anche magici. Max 6 persone al giorno. 0.5 ore di preparazione per persona.
\item \textbf{Energia/B}: Energia Positiva, Fuoco
\item \textbf{Vantaggio}: Riduzione 1 al danno non magico (RD 1-magico)
\item \textbf{Liste Magia Privilegiate}: Ammaliamento, Divinazione
\item \textbf{Arma Preferita}: Spada corta
\item \textbf{Regola}: Non rifiutare un buon bicchiere di vino
\end{itemize}

\subsubsection{Orudjs}\index{Orudjs}\label{orudjs}\index{Impulsivo}\hypertarget{orudjs}{}

\begin{enfasi}{
Nulla è più facile che illudersi. Perché l'uomo crede vero ciò che desidera. (Demostene)
}\end{enfasi}

Orudjs striscia nelle profondità delle caverne, circondato da gemme, tesori, servitori zelanti.

Descritto come una melma informe da chi ha percepito una parvenza della forma, Orudjs è Patrono della illusione e della finzione.

Con il solo pensiero convince chiunque di qualsiasi cosa voglia. Adora il teatro per ciò che per lui è, la rappresentazione della falsità, l'essere tante persone ed in realtà nessuna.

Dove domina regna il caos dove ognuno è convinto di essere nel giusto e guerre tra clan nutrono la sua fame senza fine.

Finge di ascoltare chi gli sta vicino ma in realtà non è interessato alle storie altrui perché le sue sono sempre le migliori. E' un codardo senza limiti ed un bugiardo con sempre un tornaconto.

I suoi Devoti sono creature deboli, che hanno bisogno di un padrone, di una voce che gli dica costantemente di cosa hanno bisogno e cosa vogliono.

Ma anche abili attori ed intrattenitori, spie sotto copertura, diplomatici o politicanti.

\begin{itemize}[leftmargin=*] \setlength{\itemsep}{0pt}
\item \textbf{Simbolo}: Una maschera teatrale bianca con solo la bocca aperta e gli occhi chiusi
\item \textbf{Caratteristica}: Carisma
\item \textbf{Tratti}: Impulsivo, Dissoluto, Ambizioso, Indeciso, Crudele, Compassionevole, Disonesto
\item \textbf{Manifestazione}: il suono di una risata profonda e contagiosa
\item \textbf{Somma dei Tratti in comune a 5 punti} punti: Il tuo eloquio è già leggendario. +2 alle prove di Intrattenere.
\item \textbf{Somma dei Tratti in comune a 10 punti}: Puoi lanciare Immagine Silenziosa 3 volte al giorno.
\item \textbf{Somma dei Tratti in comune a 15 punti}: Il tuo eloquio è già leggendario. +1d6 aggiuntivo alle prove di Intrattenere. Puoi lanciare \hyperlink{Immagine Maggiore}{Immagine Maggiore} 1 volta al giorno.
\item \textbf{Somma dei Tratti in comune a 20 punti}: La tua voce è suadente. Una creatura, da te individuata, che ti ascolti per almeno un minuto deve fare un Tiro Salvezza Volontà DC 30 oppure essere sotto l'influenza di Dominare Persone/Mostri. Una volta al giorno
\item \textbf{Energia/N}: Elettricità, Fuoco
\item \textbf{Vantaggio}: +4 alle prove di Ingannare
\item \textbf{Liste Magia Privilegiate}: Ammaliamento, Illusione
\item \textbf{Arma Preferita}: Stocco
\item \textbf{Regola}: Devi avere sempre l'ultima parola
\end{itemize}

\subsubsection{Orlaith}\index{Orlaith}\label{orlaith}\index{Giustizia}\hypertarget{orlaith}{}

\begin{enfasi}{
Io sono fatto per combattere il crimine, non per governarlo. Non è ancora giunto il tempo in cui gli uomini onesti possono servire impunemente la patria. I difensori della libertà saranno sempre dei proscritti finché la masnada dei furfanti dominerà. (Maximilien de Robespierre)
}\end{enfasi}

Ovvero il Patrono della Giustizia e della Vendetta. Lui segue pedissequamente le leggi e pretende che i suoi sottoposti eseguano senza alcuna discussione gli ordini impartiti. E' mosso da uno spirito gentile e buono che però tiene ben nascosto dietro le sue azioni dirette ed incisive, spudorate e mortali. Orlaith è vendetta che si fa legge. Agisce per senso di giustizia con i suoi metodi. Di lui attirano il portamento e lo sguardo fiero.

I Devoti di Orlaith spesso sono giudici e giustizieri, persone che hanno deciso di portare la giustizia ovunque, perché Orlaith non può stare fermo, c'è sempre qualcuno da giudicare e punire.

Attenti ai Seguaci di Orlaith, vanità, vendetta ed intransigenza li rendono odiosi e mal disposti verso tutti.

\begin{itemize}[leftmargin=*] \setlength{\itemsep}{0pt}
\item \textbf{Simbolo}: Una mano stesa su un libro chiuso
\item \textbf{Caratteristica}: Forza
\item \textbf{Tratti}: Vanitoso, Intransigente, Coraggioso, Testardo, Dissoluto, Vendicativo, Curioso
\item \textbf{Manifestazione}: l'immagine di una stadera, sbilanciata.
\item \textbf{Somma dei Tratti in comune a 5 punti} punti: Richiami a te 1 \hyperlink{Mastino}{mastino} che obbedisce ai tuoi comandi. Durata 1 minuto. Una volta al giorno. Costo 2 Azioni.
\item \textbf{Somma dei Tratti in comune a 10 punti}: Un paio di manette si manifesta attorno ai polsi della creatura (massimo taglia grande) entro 27 metri. Tiro Salvezza Riflessi DC 25 per annullare. Costo 2 Azioni. Una volta al giorno. Forza/Artista della Fuga DC 20 per liberarsi.
\item \textbf{Somma dei Tratti in comune a 15 punti}: Crei un raggio di Luce lungo 27 metri e largo pochi centimetri. Ogni creatura attraversata subisce 8d6 di danno, DC 25 Riflessi per dimezzare. Una volta al giorno. Costo 2 Azioni.
\item \textbf{Somma dei Tratti in comune a 20 punti}: Il tuo udito è solo per la verità. Intorno a te per 3 metri, te compreso, è sempre attiva \hyperlink{Zona di Verità}{Zona di Verità}.
\item \textbf{Energia/B}: Luce, Suono
\item \textbf{Vantaggio}: Qualsiasi arma non improvvisata nelle tue mani fa almeno 1d6 di danno
\item \textbf{Liste Magia Privilegiate}: Illusione, Fuoco
\item \textbf{Arma Preferita}: Lancia da fante
\item \textbf{Regola}: Non rifiutare un ordine da una legittima autorità
\end{itemize}

\subsubsection{Rezh}\index{Rezh}\label{rezh}\index{Avarizia}\hypertarget{rezh}{}

\begin{enfasi}{
L'avidità, non trovo una parola migliore, è valida, l'avidità è giusta, l'avidità funziona, l'avidità chiarifica, penetra e cattura l'essenza dello spirito evolutivo. L'avidità in tutte le sue forme: l'avidità di vita, di amore, di sapere, di denaro, ha improntato lo slancio in avanti di tutta l'umanità. E l'avidità, ascoltatemi bene, non salverà solamente la Teldar Carta, ma anche l'altra disfunzionante società che ha nome America. (Gordon Gekko dal film Wall Street, 1987)
}\end{enfasi}

Il Patrono che disprezza tutto. Rezh ama, vuole, tocca, rimira solo le sue monete lucide e brillanti. Non sono mai abbastanza, nessuna ricchezza è mai sufficiente per lei. Rezh, l'avara tiene tutto per se, non conosce compassione, non conosce carità, non conosce condivisione. La sua fame di denaro, di ricchezze la rende prona a qualsiasi bassezza. Disprezza tutto e tutti e giudica tutto e tutti seguendo solo il suo personale metro di giudizio. In ogni moneta c'è un pò di Rezh. Nella ossidatura di ogni moneta si può vedere l'impronta di Rezh.

Nelle profondità delle caverne i devoti di Rezh scavano cercando tesori, profanano catacombe e cacciano, cacciano ed uccidono chiunque abbia con se qualcosa di prezioso, anche il migliore amico.

Rezh si è preoccupata di distruggere in verdi fiamme qualsiasi documento finanziario potesse esistere prima della venuta dei Patroni. In un gesto che può sembrare di estrema generosità ha cancellato tutti i debite delle persone.

Tra gli umani i Devoti di Rezh diventano esploratori, tombaroli, persone sempre alla ricerca di un tesoro e di una moneta in più.

\begin{itemize}[leftmargin=*] \setlength{\itemsep}{0pt}
\item \textbf{Simbolo}: una pila di monete con un ratto vicino
\item \textbf{Caratteristica}: Intelligenza
\item \textbf{Tratti}: Avaro, Indeciso, Ambizioso, Invidioso, Crudele, Cinico, Paziente
\item \textbf{Manifestazione}: un rumore di monete che cadono avvolge l'incantatore
\item \textbf{Somma dei Tratti in comune a 5 punti} punti: Sei un esperto di monete e gemme, nessun falsario può ingannarti. +1d6 alle prove di Consapevolezza e Conoscenza relative.
\item \textbf{Somma dei Tratti in comune a 10 punti}: Usi le gemme come ricettacoli. Puoi scaricare un incantesimo di 3 livello o inferiore in una gemma, che deve avere valore minimo di 10mo x livello dell'incantesimo. La gemma conserva l'incantesimo per 6 ore. Per attivare la gemma usi 2 azioni e viene eseguito l'incantesimo che contiene consumando la gemma.
\item \textbf{Somma dei Tratti in comune a 15 punti}: Puoi tirare fuori dalle tasche 1 moneta d'oro ogni volta che vuoi. Max 10 mo al giorno. Costo 1 Azione.
\item \textbf{Somma dei Tratti in comune a 20 punti}: La tua armatura viene coperta da scintillio dorato e di gemme. Guadagni +4 alla Difesa e +2d6 Tiro Salvezza Tempra per 1 ora. Costo 2 Azioni, una volta al giorno.
\item \textbf{Energia/M}: Vuoto, Elettricità
\item \textbf{Vantaggio}: quando trovi un tesoro aumentane il valore del 1\%
\item \textbf{Liste Magia Privilegiate}: Abiurazione
\item \textbf{Arma Preferita}: Falcetto
\item \textbf{Regola}: Non lasciare un tesoro incustodito
\end{itemize}

\subsubsection{Shayalia}\index{Shayalia}\label{shayalia}\index{Lussuria}\hypertarget{shayalia}{}

\begin{enfasi}{
Chi pianta un giardino semina felicità (proverbio cinese)

\medskip

Nulla di grande al mondo è stato compiuto senza passione. (Georg Wilhelm Friedrich Hegel)

}\end{enfasi}

Patrono dell'Arcano di Tenebra. Shayalia è l'anima oscura della perdizione, del tradimento, della lussuria più sordida e peccaminosa. Adora i bordelli. Le piace l'odore acre del sudore, la pelle lucida di oli e profumi. Le passioni, le vendette che li si consumano, la distruzione fisica e morale che in quei luoghi viene perpetrata è la sua vita.

Shayalia è una donna, spesso sola, che gode e vive i piaceri più fisici. Vive di vendette lungamente e ben dettagliatamente progettate. Vendicativa ed amorale, non giudica con metro di giudizio umano, il suo godere non è neppure lontanamente comprensibile. Shayalia è quanto di più vicino a Calicante sia stato creato. Sono le passioni, le pulsioni, i liquidi umorali che la fanno inebriare.

Shayalia è la concubina che ti ammalia e ti distrugge, goccia dopo goccia. I veleni sono le sue armi, le debolezze umane il suo campo.

I Devoti di Shayalia sono spie, figli bastardi, amanti di potenti signori che agiscono all'ombra.

Ljust disgustata dalla visione di un Patrono del genere instillò in Shayalia l'amore e passione per le piante ed animali. E così molti dei più famosi botanici, erboristi e zoologi sono Devoti di Shayalia, forse le uniche cose che Shayalia veramente può amare.

\medskip

\begin{narratore}[Efrem e Shayalia]
Mentre \textbf{Efrem} è il patrono della natura selvaggia, Shayalia incarna la devozione e l'amore umano per la natura. Il primo sovraintende dall'alto la natura incontaminata, la seconda si cala e diventa tutt'uno con essa costruendo magnifici giardini.

\medskip

\textbf{Sumkjr} e \textbf{Shayalia} sono complementari nel tenere in mano gli umori sfuggenti delle creature. Agiscono come espressione diretta dei Patroni della Genesi.

\end{narratore}

\medskip

\begin{itemize}[leftmargin=*] \setlength{\itemsep}{0pt}
\item \textbf{Simbolo}: un cuscino stropicciato sporco di sangue
\item \textbf{Caratteristica}: Carisma
\item \textbf{Tratti}: Dissoluto, Cinico, Crudele, Vendicativo, Paziente, Compassionevole, Vanitoso
\item \textbf{Manifestazione}: il Devoto è avvolto da un mantello di velluto nero
\item \textbf{Somma dei Tratti in comune a 5 punti} punti: I tempi per preparare una pozione sono dimezzati. Gli incantesimi di Cura hanno effetto anche su animali e piante.
\item \textbf{Somma dei Tratti in comune a 10 punti}: Il tuo tocco è vita per la natura. I tuoi incantesimi di cura agiscono su animali e piante naturali in maniera massimizzata.
\item \textbf{Somma dei Tratti in comune a 15 punti}: Dal tuo palmo secerni veleno. Il tuo tocco, o tramite arma in mischia veicola il veleno. Tiro Salvezza Tempra DC 25 o -2 a Saggezza e Destrezza per 10 minuti, un soggetto avvelenato non può esserlo nuovamente per 24 ore. Costo 1 Azione.
\item \textbf{Somma dei Tratti in comune a 20 punti}: Il tuo tocco è vita per la natura. Puoi curare animali e piante magiche. Sei immune ai veleni naturali. +1d6 Conoscenza Natura.
\item \textbf{Energia/M}: Vuoto, Elettricità
\item \textbf{Vantaggio}: +4 TS contro Veleni
\item \textbf{Liste Magia Privilegiate}: Illusione oppure Animali e Piante ed una Liste Magia Elementale
\item \textbf{Arma Preferita}: Frusta
\item \textbf{Regola}: Non rinunciare ad umiliare
\end{itemize}

\subsubsection{Sixiser}\index{Sixiser}\label{sixiser}\index{Ingordigia}\hypertarget{sixiser}{}

\begin{enfasi}{
La forza che si oppone al destino è in realtà una debolezza. (Franz Kafka)

\medskip

Ora stai zitto, bimbo, non piangere,

mamma realizzerà tutti i tuoi incubi. (Mother, 1979 The Wall, Pink Floyd)

}\end{enfasi}

Il Patrono che è indifferente al presente in quanto totalmente, compulsivamente ossessionato dal futuro e dal suo destino. Negli angoli più remoti dei mondi conosciuti si narra che Sixiser accumuli di tutto, indifferente a tutto e tutti.

Terrorizzato dal futuro che vede, da una ipotetica fine di sé e del tutto vive una vita di ritiro, spirituale e fisico. Si priva volontariamente di tutto il necessario. Ma allo stesso accumula qualunque oggetto incroci la sua strada nella speranza di un ritorno.

E' paranoico e non si fida di nessuno. Usa i suoi poteri di divinazione per conoscere e scrutare tutti.

Sixiser è padrone degli incubi, dei sogni più spaventosi delle visioni di morte. Spesso usa gli incubi quale mezzo di comunicazione con i suoi seguaci.

I Devoti di Sixiser sono spesso negromanti circondati da non morti ed altre creature silenziose ed ubbidienti. Chi si rifugia alla ricerca della solitudine e dello studio, chi mira ad espandere e governare intere città e nazioni al fine di sentirsi più sicuro, è Devoto a Sixiser.

\begin{itemize}[leftmargin=*] \setlength{\itemsep}{0pt}
\item \textbf{Simbolo}: Un forziere straripante di ogni cosa che non si può chiudere
\item \textbf{Caratteristica}: Saggezza
\item \textbf{Tratti}: Prudente, Indeciso, Intransigente, Impulsivo, Disonesto, Cinico, Sospettoso
\item \textbf{Manifestazione}: due mani che circondano, come a nascondere, la testa dell'incantatore
\item \textbf{Somma dei Tratti in comune a 5 punti} punti: acquisisci la visione crepuscolare fino 9 metri, o 18 metri se già presente.
\item \textbf{Somma dei Tratti in comune a 10 punti}: vedi nell'oscurità anche magica entro 9 metri. Individui automaticamente le trappole non magiche entro 3 metri da te.
\item \textbf{Somma dei Tratti in comune a 15 punti}: Toccando un oggetto sei in grado di capirne tutte le proprietà magiche e non, anche se è maledetto. 3 volte al giorno.
\item \textbf{Somma dei Tratti in comune a 20 punti}: Sei in grado di animare una creatura morta da non più di un giorno come non morto da 1 grado di Sfida (tipo zombi/scheletro a secondo dello stato). Una volta al giorno. Costo 2 Azioni.
\item \textbf{Energia/M}: Elettricità, Energia Negativa
\item \textbf{Vantaggio}: immune alle malattie naturali
\item \textbf{Liste Magia Privilegiate}: Necromanzia
\item \textbf{Arma Preferita}: Falcione
\item \textbf{Regola}: Non fidarti
\end{itemize}

\subsubsection{Sumkjr}\index{Sumkjr}\label{sumkjr}\index{Gentilezza}\hypertarget{sumkjr}{}

\begin{enfasi}{
Ogni uomo è colpevole di tutto il bene che non ha fatto. (Voltaire)

\medskip

Tutto ciò che non viene donato va perduto. (Dominique Lapierre)

}\end{enfasi}

Patrono dell'Arcano di Luce. Sumkjr è bontà, correttezza, lealtà, giustizia, protezione.

Sumkjr è il cavaliere che protegge gli innocenti, è la spada di Ljust nella battaglia finale. Difende i deboli e lenisce le ferite.

Sumkjr porta la Luce di Ljust ovunque, nessun pericolo potrà mai fermare Sumkjr dalla sua continua, infinita, cerca del bene.

Un Devoto di Sumkjr agisce lealmente e con onore, sempre perseguendo il bene ultimo, il suo essere non può essere piegato al male, all'ingiustizia, al disonore.

Con coraggio e determinazione il Devoto affronta ogni sfida ma non solo per senso del dovere, ma perché profondamente votato al suo destino. Sumkjr sa che poche persone reggono tale standard perché a differenza dei Devoti della Patrona delle Genesi i suoi Devoti non nascono per essere tali, ma lo diventano grazie alla loro profonda e determinata forza di volontà.

Per questo motivo Ljust interviene in loro favore con l'elaborato Rito del Rinnovo, grazie al quale ogni anno al Devoto meritevole viene fatto aggiungere un punto Tratto.

Sumkjr è un soldato valoroso, il migliore amico del giusto.

Calicante, preso dall'orrore alla vista di un Patrono così fatto, lo privò della capacità di amare e provare veri sentimenti d'affetto. Portare il bene per un Devoto di Sumkjr è un qualcosa di normale come è normale non riuscire ad essere empatico con chi soffre. Il Devoto sa cosa deve fare e perché, ma non riesce a commuoversi od amare di fronte alle sofferenze od alle carezze di una donna/uomo.

\begin{itemize}[leftmargin=*] \setlength{\itemsep}{0pt}
\item \textbf{Simbolo}: tre gocce di sangue che cadono una dietro l'altra
\item \textbf{Caratteristica}: Carisma
\item \textbf{Tratti}: Gentile, Coraggioso, Testardo, Sospettoso, Altruista, Curioso, Estroverso
\item \textbf{Manifestazione}: il Devoto è avvolto da un mantello di broccato dorato
\item \textbf{Somma dei Tratti in comune a 2 punti} punti: Il tocco della tua spada é vita. Una creatura toccata con la tua arma recupera 3d6 Punti Ferita. Una volta al giorno. Costo 2 Azioni.
\item \textbf{Somma dei Tratti in comune a 7 punti}: La tua Volontà é più forte del metallo. Guadagni un +2 ai Tiri Salvezza su Volontà
\item \textbf{Somma dei Tratti in comune a 11 punti}: Puoi lanciare l'incantesimo Cono di Freddo, 40 danni da Elettricità. DC 25 per dimezzare. Una volta al giorno. Costo 2 Azioni.
\item \textbf{Somma dei Tratti in comune a 15 punti}: Sacrifichi la tua vita per portare in vita una creatura morta da non più di 1 settimana. Una volta. Costo 3 Azioni.\index{Resuscitare}\\
\textbf{Nota}: Sumkjr è l'unico Patrono a concedere benefici per una somma Tratti minore del normale
\item \textbf{Energia/B}: Energia Positiva, Elettricità
\item \textbf{Vantaggio}: +1 a tutti i Tiri Salvezza
\item \textbf{Liste Magia Privilegiate}: Cura
\item \textbf{Arma Preferita}: Spada Bastarda
\item \textbf{Regola}: Non compiere atti sessuali
\end{itemize}

\textbf{Le 7 Regole Luminose}\index{7 Regole Luminose}

Le Sette regole Luminose sono un insieme di norme e comportamenti tenuti, a vario titolo, dai Devoti che vogliono seguire la strada della Luce di Ljust.

I Devoti di Sumkjr devono seguirle tutte e 7, altri Devoti di altri Patroni, sempre positivi od almeno neutrali, seguono solo alcune di questi dettami, come regola per non cadere nelle braccia di Calicante.

\begin{enumerate}[leftmargin=*] \setlength{\itemsep}{0pt}
\item Proteggi i deboli e chi non sa difendersi dai soprusi
\item Ama la vita e proteggila
\item Combatti contro le ingiustizie e chi porta sofferenze e dolore
\item Lenisci le ferite ed i dolori. Placa gli animi e favorisci la pace ed armonia
\item Onestà e Lealtà sono le tua fondamenta
\item Sei un maestro di virtù. Fa che gli altri possano prendere ispirazione dalle tue gesta
\item Non lasciare che la tua inazione generi sofferenza
\end{enumerate}

\subsubsection{Tàhil}\index{Tàhil}\label{tahil}\index{Vendicativo}\hypertarget{tahil}{}

\begin{enfasi}{
Le persone malvagie fanno cose malvagie perché possono. (Sherlock Holmes - Gioco di ombre)
}\end{enfasi}

Il Patrono che brama solo infliggere dolore. Tàhil è l'incarnazione della furia di vendetta, della volontà di fare del male per il piacere della sofferenza.

Tàhil è stato creato quando Calicante ha scoperto la morte del suo primo figlio, è stata una creazione di puro istinto e furia, di rabbia e desiderio di distruzione totale, di sofferenza ed umiliazione; nulla di meglio che creare il Drago.

Tàhil ha ben pochi Devoti o Seguaci che si dichiarino tali, ma a differenza di tutti gli altri Patroni a Tàhil la cosa non interessa, lui è la manifestazione dell'odio puro e questo gli basta.

\begin{itemize}[leftmargin=*] \setlength{\itemsep}{0pt}
\item \textbf{Simbolo}: Una svastica
\item \textbf{Caratteristica}: Forza
\item \textbf{Tratti}: Vendicativo, Disonesto, Arrogante, Cinico, Ambizioso, Testardo, Impulsivo
\item \textbf{Manifestazione}: un rumore di tuono
\item \textbf{Somma dei Tratti in comune a 5 punti}: Puoi aggiungere 1d6 ad un Tiro per Colpire. Una volta al giorno, prima di effettuare il Tiro per Colpire, come Azione Immediata.
\item \textbf{Somma dei Tratti in comune a 10 punti}: Tre volte al giorno, prima di effettuare il Tiro per Colpire, puoi dichiarare di colpire. Azione Immediata
\item \textbf{Somma dei Tratti in comune a 15 punti}: Un tuo colpo andato a segno causa almeno un Tiro Critico.
\item \textbf{Somma dei Tratti in comune a 20 punti}: Il dado del Critico dell'arma aumenta di una taglia.
\item \textbf{Energia/M}: Fuoco, Vuoto
\item \textbf{Vantaggio}: Riduzione danno da Luce 10
\item \textbf{Liste Magia Privilegiate}: Ammaliamento
\item \textbf{Arma Preferita}: Spadone
\item \textbf{Regola}: Mai essere clemente
\end{itemize}

\subsubsection{Tazher}\index{Tazher}\label{tazher}\index{Malizia}\hypertarget{tazher}{}

\begin{enfasi}{
Una persona spesso finisce con l'assomigliare alla sua ombra. (Rudyard Kipling)
}\end{enfasi}

Il Patrono delle Ombre; colui che silenzioso, ti uccide. Non saprai mai il perché. Non conoscerai mai il suo aspetto ma, se improvvisamente hai una sensazione di gelo, Tazher è dietro di te pronto a prendere la tua vita.

Doppiogiochista dall'animo cattivo, chiedi il suo aiuto solo se sei disposto a pagarne il prezzo che lui e lui solo deciderà.

Vive di oscurità e sangue. Le ombre sono le sue amiche e la tenebra il suo mantello.

Si circonda di assassini, mercenari, chiunque uccida senza provare sentimenti. Nelle profondità del sottosuolo nutre i suoi adepti di dolore, sangue e morte.

Ljust inorridita da tanto odio e nichilismo instillò nel Patrono il rispetto per i morti. Un Devono si Tazher non si accanirà contro un defunto ne violerà il suo cadavere. Molti cacciatori di non morti sono devoti di Tazher.

L'umano Devoto di Tazher è il ladro, l'assassino, il bandito, chiunque viva per l'oscurità ed il proprio tornaconto. Un Devoto di Tazher è estremamente pericoloso in combattimento.

\begin{itemize}[leftmargin=*] \setlength{\itemsep}{0pt}
\item \textbf{Simbolo}: Lo scintillio della lama nel buio
\item \textbf{Caratteristica}: Destrezza
\item \textbf{Tratti}: Disonesto, Ambizioso, Paziente, Cinico, Indeciso, Arrogante, Crudele
\item \textbf{Manifestazione}: l'ombra del Devoto prende vita muovendo l'arma
\item \textbf{Somma dei Tratti in comune a 5 punti} punti: Guadagni +2 alle prove di Furtività.
\item \textbf{Somma dei Tratti in comune a 10 punti}: La tua Scurovisione diventa di 6 metri.
\item \textbf{Somma dei Tratti in comune a 15 punti}: finché cammini sopra delle ombre o al buio (oscurità) sei invisibile. Puoi essere comunque rilevato con la luce o incantesimi di divinazione.
\item \textbf{Somma dei Tratti in comune a 20 punti}: L'oscurità non è più un problema. Vedi nell'oscurità anche magica come se fosse giorno. Quando sei in un ambiente in piena luce sei abbagliato con un -2 al Tiro per Colpire.
\item \textbf{Energia/M}: Vuoto, Ghiaccio
\item \textbf{Vantaggio}: Scurovisione 3 metri
\item \textbf{Liste Magia Privilegiate}: Trasmutazione
\item \textbf{Arma Preferita}: Falcione in asta
\item \textbf{Regola}: 5 Secondi. Il tempo per rubare ad un morto, non di più.
\end{itemize}

\subsubsection{Thaft}\index{Thaft}\label{thaft}\index{Pazienza}\hypertarget{thaft}{}

\begin{enfasi}{
La morte, il più atroce dunque di tutti i mali, non esiste per noi. Quando noi viviamo la morte non c'è, quando c'è lei non ci siamo noi. Non è nulla né per i vivi né per i morti. Per i vivi non c'è, i morti non sono più. (Epicuro)
}\end{enfasi}

Il Patrono che accompagna nella nascita e nella morte. Silenzioso, resta in disparte e osserva lo scorrere della vita degli uomini. Quasi umile nella sua semplicità, Thaft è ovunque. Testimone silenzioso della vita umana; nel momento in cui una vita scivola via, Thaft assiste, nell'attimo in cui una vita nasce, Thaft è presente.

Thaft sa anche che non si può essere sempre e solo osservatori. Attraverso il suo taccuino sacro e magico può decidere e giudicare della vita degli uomini, perché se una spada ferisce, è solo Thaft che ne decide la morte.

I Devoti di Thaft sono i sacerdoti dell'ultimo viaggio, coloro che proteggono e vegliano sulle anime e corpi dei morti. Profondamente contrari all'utilizzo dei non-morti ne perseguono la distruzione.

Un Devoto di Thaft rispetta la vita come la morte e non teme di arrecare distruzione per un equilibrio maggiore.

Thaft è stato plasmato da Atmos.

\begin{itemize}[leftmargin=*] \setlength{\itemsep}{0pt}
\item \textbf{Simbolo}: Un libro aperto con un teschio sopra
\item \textbf{Caratteristica}: Saggezza
\item \textbf{Tratti}: Codardo, Paziente, Estroverso, Leale, Gentile, Vanitoso, Vendicativo
\item \textbf{Manifestazione}: si sente il pianto di un bambino appena nato o il sospiro della morte
\item \textbf{Somma dei Tratti in comune a 5 punti} punti: Il tuo tocco è letale per i non morti. Un tuo tocco infligge 2d6 di danno ad un non morto. Costo 2 Azioni compreso il tocco. Fino a 3 volte al giorno.
\item \textbf{Somma dei Tratti in comune a 10 punti}: Il tuo tocco lenisce. Una volta al giorno puoi rimuovere Cecità o Sordità. Costo 2 Azioni.
\item \textbf{Somma dei Tratti in comune a 15 punti}: Un non morto, con GS inferiore alla somma dei tuoi Tratti in comune, deve effettuare un Tiro Salvezza Tempra DC 30 o essere distrutto se toccato dalla tua mano. Costo 2 Azioni.
\item \textbf{Somma dei Tratti in comune a 20 punti}: Uccidi la creatura toccata. Tiro Salvezza su Volontà DC 30 o morte. Una volta alla settimana. Costo 2 Azioni.
\item \textbf{Energia/N}: Suono, Elettricità
\item \textbf{Vantaggio}: Riduzione danno 5 a Vuoto e Luce
\item \textbf{Liste Magia Privilegiate}: Necromanzia, Animali e piante
\item \textbf{Arma Preferita}: Arco
\item \textbf{Regola}: Non creare un non morto
\end{itemize}

\subsubsection{Torbiorn}\index{Torbiorn}\label{torbion}\index{Crudele}\hypertarget{torbiorn}{}

\begin{enfasi}{
Riguardo all'arroganza, i violenti la soffrono, ma i saggi la deridono. (Tito Livio, attribuita ad Astimede)
}\end{enfasi}

Il Patrono che meglio incarna il concetto \emph{non è mai abbastanza}. Alto, bello come una statua classica ma, proprio come quest'ultima, senza calore e vita, Torbiorn rasenta la perfezione maniacale nel vestirsi, nell'atteggiarsi.

Nulla è mai abbastanza per lui. Nessuno è mai alla sua altezza. Ed eccolo che con arroganza e ironia va a modificare tutto il modificabile per poter placare questa profonda insoddisfazione. Qualora il risultato finale raggiunto non lo soddisfi, e accade molto spesso, ecco che prende il sopravvento il suo cinismo e distrugge tutto senza curarsi della sofferenza che sta arrecando a chi gli sta attorno.

Nelle lande abbandonate il Devoto di Torbiorn è il Tiranno dal pugno di ferro che agisce solo per proprio desiderio e piacere senza curarsi di nessun'altro.

Il Devoto di Torbiorn è il tipico aristocratico ricco e svogliato,colui che cerca sempre la strada più facile e meno rischiosa.

Incurante degli altri si diverte nello sfruttare i lavori altrui e trarne giovamento.

\begin{itemize}[leftmargin=*] \setlength{\itemsep}{0pt}
\item \textbf{Simbolo}: Uno specchio opaco
\item \textbf{Caratteristica}: Carisma
\item \textbf{Tratti}: Crudele, Impulsivo, Arrogante, Disonesto, Cinico, Indeciso, Compassionevole
\item \textbf{Manifestazione}: schegge di specchio rotto tutto intorno al Devoto come un turbine
\item \textbf{Somma dei Tratti in comune a 5 punti} punti: Con un gesto puoi rinfrescare i tuoi vestiti e te stesso rendendoli puliti e profumati. Costo 1 Azione. 3 volte al giorno.
\item \textbf{Somma dei Tratti in comune a 10 punti}: Il tuo sputo è velenoso. Se il Tiro per Colpire a tocco va a segno -2 Forza, non cumulabile. Durata 1 minuto. Tre volte al giorno. Costo 1 Azione.
\item \textbf{Somma dei Tratti in comune a 15 punti}: Fissando l'obiettivo negli occhi lo costringi a fermarsi. Il soggetto non può eseguire Azioni di Movimento. Tiro Salvezza su Volontà DC 30. Una volta al giorno. Costo 2 Azioni.
\item \textbf{Somma dei Tratti in comune a 20 punti}: Dalle tue dita partono dei viticci che pungono fino a 10 avversari. Ogni viticcio, lungo fino a 18 metri causa 2d6 di danno, Tiro Salvezza Riflessi DC 25 per dimezzare. Costo 2 Azioni.
\item \textbf{Energia/N}: Fuoco, Suono
\item \textbf{Vantaggio}: +3 contro gli incantesimi della Lista Divinazione.
\item \textbf{Liste Magia Privilegiate}: Trasmutazione
\item \textbf{Arma Preferita}: Ascia ad una mano
\item \textbf{Regola}: Non essere sciatto, mal vestito o disordinato.
\end{itemize}

\medskip

\begin{narratore}[Adeguare ed adeguarsi]

In accordo con il Narratore, ed adeguatamente motivato, è possibile cambiare Vantaggio e Liste di Magia Privilegiate.

\end{narratore}

\subsubsection{Elenco Patrono - Tratto}\index{Elenco Patrono - Tratto}\hypertarget{tabellacollegamentopatronotratto}{}\label{tabellacollegamentopatronotratto}

\smallskip

I Patroni sono indicati in ordine alfabetico per Tratto maggiormente caratterizzante.

\medskip

\setlength{\parindent}{0cm}{

\textbf{\hyperlink{gaya}{Gaya}}: Altruista, Gentile, Sospettoso, Cinico, Invidioso, Disonesto, Arrogante

\smallskip

\textbf{\hyperlink{calicante}{Calicante}}: Ambizioso, Disonesto, Vendicativo, Cinico, Dissoluto, Arrogante, Avaro

\smallskip

\textbf{\hyperlink{erondil}{Erondil}}: Arrogante, Vendicativo, Ambizioso, Compassionevole, Entusiasta, Leale, Avaro

\smallskip

\textbf{\hyperlink{rezh}{Rezh}}: Avaro, Indeciso, Ambizioso, Invidioso, Crudele, Cinico, Paziente

\smallskip

\textbf{\hyperlink{sixiser}{Sixiser}}: Prudente, Indeciso, Intransigente, Impulsivo, Disonesto, Cinico, Sospettoso

\smallskip

\textbf{\hyperlink{cattalm}{Cattalm}}: Cinico, Arrogante, Ambizioso, Intransigente, Dissoluto, Sospettoso, Paziente

\smallskip

\textbf{\hyperlink{ljust}{Ljust}}: Compassionevole, Testardo, Coraggioso, Estroverso, Altruista, Leale, Paziente

\smallskip

\textbf{\hyperlink{gradh}{Gradh}}: Coraggioso, Vanitoso, Arrogante, Gentile, Invidioso, Leale, Sospettoso

\smallskip

\textbf{\hyperlink{nihar}{Nihar}}: Curioso, Coraggioso, Compassionevole, Vanitoso, Invidioso, Avaro, Crudele

\smallskip

\textbf{\hyperlink{tazher}{Tazher}}: Disonesto, Ambizioso, Paziente, Cinico, Indeciso, Arrogante, Crudele

\smallskip

\textbf{\hyperlink{shayalia}{Shayalia}}: Dissoluto, Cinico, Crudele, Vendicativo, Paziente, Compassionevole, Vanitoso

\smallskip

\textbf{\hyperlink{ledyal}{Ledyal}}: Entusiasta, Compassionevole, Prudente, Gentile, Curioso, Codardo, Testardo

\smallskip

\textbf{\hyperlink{nethergal}{Nethergal}}: Estroverso, Curioso, Testardo, Vanitoso, Vendicativo, Arrogante, Paziente

\smallskip

\textbf{\hyperlink{sumkjr}{Sumkjr}}: Gentile, Coraggioso, Testardo, Sospettoso, Altruista, Curioso, Estroverso

\smallskip

\textbf{\hyperlink{atmos}{Atmos}}: Indeciso, Prudente, Intransigente, Paziente, Vendicativo, Curioso, Avaro

\smallskip

\textbf{\hyperlink{krondal}{Krondal}}: Intransigente, Vanitoso, Arrogante, Sospettoso, Paziente, Ambizioso, Testardo

\smallskip

\textbf{\hyperlink{belevon}{Belevon}}: Invidioso, Ambizioso, Dissoluto, Disonesto, Compassionevole, Paziente, Altruista

\smallskip

\textbf{\hyperlink{orudjs}{Orudjs}}: Impulsivo, Dissoluto, Ambizioso, Indeciso, Crudele, Compassionevole, Disonesto

\smallskip

\textbf{\hyperlink{efrem}{Efrem}}: Leale, Indeciso, Prudente, Impulsivo, Testardo, Paziente, Ambizioso

\smallskip

\textbf{\hyperlink{torbiorn}{Torbiorn}}: Crudele, Impulsivo, Arrogante, Disonesto, Cinico, Indeciso, Compassionevole

\smallskip

\textbf{\hyperlink{nedraf}{Nedraf}}: Paziente, Vanitoso, Coraggioso, Intransigente, Entusiasta, Arrogante, Cinico

\smallskip

\textbf{\hyperlink{atherim}{Atherim}}: Sospettoso, Compassionevole, Altruista, Intransigente, Coraggioso, Entusiasta, Vanitoso\

\smallskip

\textbf{\hyperlink{thaft}{Thaft}}: Codardo, Paziente, Estroverso, Leale, Gentile, Vanitoso, Vendicativo

\smallskip

\textbf{\hyperlink{lynx}{Lynx}}: Testardo, Coraggioso, Cinico, Intransigente, Vendicativo, Estroverso, Vanitoso

\smallskip

\textbf{\hyperlink{orlaith}{Orlaith}}: Vanitoso, Intransigente, Coraggioso, Testardo, Dissoluto, Vendicativo, Curioso

\smallskip

\textbf{\hyperlink{laydel}{Laydel}}: Vendicativo, Paziente, Ambizioso, Intransigente, Invidioso, Cinico, Arrogante

\smallskip

\textbf{\hyperlink{tahil}{Tàhil}}: Vendicativo, Disonesto, Arrogante, Cinico, Ambizioso, Testardo, Impulsivo

}

\end{multicols}

\vfill

\begin{enfasi}
	Gli dei tessono sventure per gli uomini, perché le generazioni future abbiano qualcosa da cantare.(Iliade, Omero)
\end{enfasi}

%\vfill

%\begin{center}
%\includegraphics[keepaspectratio,width=0.6\textwidth]{immagini/archetipijung.png}

%\emph{I 12 Archetipi di Jung}
%\end{center}

\pagebreak

